\documentclass[10pt]{article}
\usepackage[a4paper,margin=0.88in]{geometry}
\usepackage{amsmath,amssymb,amsthm,mathtools,amscd,booktabs,enumitem,xcolor,hyperref}
\hypersetup{colorlinks=true,linkcolor=blue!55!black,citecolor=blue!55!black,urlcolor=blue!55!black}
\setlength{\emergencystretch}{3em}
\setlength{\parindent}{0pt}
\setlength{\parskip}{0.55em}
\setlist{leftmargin=1.6em,itemsep=0.25em,topsep=0.25em}
\sloppy

\newtheorem{theorem}{Theorem}[section]
\newtheorem{lemma}[theorem]{Lemma}
\newtheorem{proposition}[theorem]{Proposition}
\newtheorem{corollary}[theorem]{Corollary}
\theoremstyle{definition}
\newtheorem{definition}[theorem]{Definition}
\newtheorem{convention}[theorem]{Convention}
\newtheorem{remark}[theorem]{Remark}

\newcommand{\Fq}{\mathbf F_q}
\newcommand{\Fbar}{\overline{\mathbf F}_q}
\newcommand{\Ql}{\mathbf Q_\ell}
\newcommand{\Qlb}{\overline{\mathbf Q}_\ell}
\newcommand{\Aone}{\mathbf A^1}
\newcommand{\Dbc}{D^b_c}
\newcommand{\Dbcint}{D^b_{c,\mathrm{int}}}
\newcommand{\Dbclat}{D^b_{c,\mathrm{lat}}}
\newcommand{\DbcW}{D^b_{c,W}}
\newcommand{\Perv}{\operatorname{Perv}}
\newcommand{\PervW}{\operatorname{Perv}_W}
\newcommand{\Rep}{\operatorname{Rep}}
\newcommand{\Sw}{\operatorname{Sw}}
\newcommand{\Frob}{\operatorname{Frob}}
\newcommand{\FT}{\operatorname{FT}}
\newcommand{\RGammac}{R\Gamma_c}
\newcommand{\deps}[1]{\par\smallskip\noindent\textit{Prerequisites.} #1}

\begin{document}
\begin{center}
{\large\bfseries Laumon's Local Fourier Theorem\par}
\medskip
{\normalsize Local transforms, stationary phase, and monodromy weights\par}
\end{center}

\begin{abstract}
This note fixes the coefficient, local-monodromy, middle-extension, and
Fourier conventions needed for Laumon's affine-line theorem.  It states the
three local Fourier transforms with their precise representation categories,
the type-3 local exact sequence, and the realness and weight results used in
the theorem.  All derived categories are ordinary Verdier triangulated
categories.
\end{abstract}

\tableofcontents

\section{Coefficients and derived categories}

Fix a finite field \(k=\Fq\) of characteristic \(p\), a prime
\(\ell\ne p\), and a finite extension \(E/\Ql\).  Enlarge \(E\), when
necessary, so that it contains the values of a fixed nontrivial additive
character
\[
                  \psi:k\longrightarrow E^\times .
\]
Fix an inclusion \(E\subset\Qlb\) and an embedding
\(\iota:\Qlb\hookrightarrow\mathbf C\).  A statement with
\(\Qlb\)-coefficients will be interpreted by the finite-stage category
defined below.

\begin{definition}[Pro-etale coefficient sheaves]
Let \(\mathcal O_E\) be the ring of integers of \(E\), choose a uniformizer
\(\varpi\), and put \(\mathcal O_{E,n}=\mathcal O_E/\varpi^n\).  For a
quasi-compact and quasi-separated scheme \(X\) of finite Krull dimension on
which \(\ell\) is invertible, let
\((\mathcal O_{E,n})_X\) be the sheafification on \(X_{\mathrm{pro\acute et}}\)
of
\[
 U\longmapsto
 \operatorname{Cont}\bigl(\pi_0(|U|),\mathcal O_{E,n}\bigr).
\]
Define the integral and rational coefficient sheaves by
\[
 (\mathcal O_E)_X=\varprojlim_n(\mathcal O_{E,n})_X,
 \qquad
 E_X=(\mathcal O_E)_X[1/\varpi]
     =\left(\varprojlim_n(\mathcal O_{E,n})_X\right)[1/\varpi].
\]
Let \(D(X,\mathcal O_E)\) and \(D(X,E)\) be the ordinary Verdier derived
categories of sheaves of modules over \((\mathcal O_E)_X\) and \(E_X\),
respectively, and define \(D(X,\mathcal O_{E,n})\) in the same way from
\((\mathcal O_{E,n})_X\).  They are obtained from cochain complexes by inverting
quasi-isomorphisms; their shifts and distinguished triangles come from
shifted complexes and mapping cones.
This scope includes schemes of finite type over \(k\) or \(\Fbar\), the
henselian traits used for local Fourier transform, their fiber products over
\(k\), and the geometric generic and special fibers used below.
\end{definition}

\begin{definition}[Derived-complete integral constructibility]
An \((\mathcal O_{E,n})_X\)-module sheaf is constructible if it is locally
constant with finite \(\mathcal O_{E,n}\)-module stalks on the members of a
finite locally closed stratification of \(X\).  Let
\(\Dbc(X,\mathcal O_{E,n})\) be the full subcategory of
\(D(X,\mathcal O_{E,n})\) with bounded constructible cohomology.

For \(K\in D(X,\mathcal O_E)\), put
\[
 K_n=K\otimes^L_{(\mathcal O_E)_X}(\mathcal O_{E,n})_X.
\]
The object \(K\) belongs to \(\Dbc(X,\mathcal O_E)\) when the following
conditions hold: every \(K_n\) is bounded constructible; one finite interval
contains the cohomological support of every \(K_n\); the canonical maps
\[
 K_n\otimes^L_{(\mathcal O_{E,n})_X}(\mathcal O_{E,m})_X
       \xrightarrow{\sim}K_m\qquad(m\le n)
\]
are isomorphisms; and the derived-completeness map
\[
                         K\longrightarrow R\!\varprojlim_nK_n
\]
is an isomorphism.  Here the derived inverse limit is taken after regarding
the reductions as a tower of \((\mathcal O_E)_X\)-module sheaves.  These
conditions define a full ordinary triangulated subcategory.
\end{definition}

\begin{definition}[Rational constructibility and lattices]
An \(E_X\)-module sheaf is lisse of finite rank if it is locally isomorphic
on the pro-etale site to \(E_X^d\) for some finite \(d\).  It is
constructible if its restriction to each member of some finite locally
closed stratification is lisse of finite rank.  Let
\(\Dbcint(X,E)\) be the full subcategory of \(D(X,E)\) whose
objects have bounded constructible cohomology.

Rationalization is the triangulated functor
\[
 \operatorname{rat}_E:D(X,\mathcal O_E)\longrightarrow D(X,E),
 \qquad K\longmapsto K\otimes^L_{(\mathcal O_E)_X}E_X.
\]
An integral lattice in \(L\in D(X,E)\) is a pair consisting of
\(K\in\Dbc(X,\mathcal O_E)\) and an isomorphism
\[
                 K\otimes^L_{(\mathcal O_E)_X}E_X\xrightarrow{\sim}L.
\]
Write \(\Dbclat(X,E)\) for the full subcategory of objects
admitting such a lattice.
\end{definition}

\begin{theorem}[Integral-lattice and intrinsic comparison]
For the schemes considered here,
\[
       \Dbclat(X,E)=\Dbcint(X,E).
\]
Their common category is denoted \(\Dbc(X,E)\).  It is an
idempotent-complete full triangulated subcategory of \(D(X,E)\), every
object admits an integral lattice, and every object has finite
Tor-amplitude.  More precisely, if \(\Dbc(X,\mathcal O_E)[1/\varpi]\) has
the same objects as \(\Dbc(X,\mathcal O_E)\) and morphism groups
\[
 \operatorname{Hom}(K,L)\otimes_{\mathcal O_E}E,
\]
then rationalization induces an exact equivalence
\[
 \Dbc(X,\mathcal O_E)[1/\varpi]^{\natural}
       \xrightarrow{\sim}\Dbc(X,E),
\]
where \((-)^{\natural}\) denotes Karoubi completion.
\end{theorem}

\begin{definition}[The \(\Qlb\)-coefficient category]
Let all finite extensions of \(\Ql\) be taken inside \(\Qlb\).  Form a
category \(\mathcal C_X\) whose objects are pairs \((E,K)\) with
\(K\in\Dbc(X,E)\), and whose morphisms are
\[
 \operatorname{Hom}_{\mathcal C_X}((E,K),(F,L))
 =\varinjlim_{M\supset E,F}
 \operatorname{Hom}_{\Dbc(X,M)}
 \bigl(K\otimes_E^LM,L\otimes_F^LM\bigr).
\]
The colimit runs over common finite extensions \(M/\Ql\), and composition
is performed after moving representatives to one common finite stage.
Define
\[
                         \Dbc(X,\Qlb)=\mathcal C_X^{\natural}.
\]
Concretely, an object is a triple \((E,K,e)\), where \(e^2=e\) is an
endomorphism of \((E,K)\) in \(\mathcal C_X\).  A morphism
\((E,K,e)\to(F,L,e')\) is a common-stage class \(u\) satisfying
\(u=e'u=ue\).  Shifts are \((E,K[1],e[1])\), and distinguished triangles
are the triangles isomorphic to direct summands of triangles represented at
a common finite extension of \(\Ql\).  This is the ordinary triangulated
Karoubi completion of the filtered finite-stage category.
\end{definition}

\begin{theorem}[Finite scalar extension]
For a finite extension \(E'/E\), derived tensor product defines an exact
functor
\[
 -\otimes_E^LE':\Dbc(X,E)\longrightarrow\Dbc(X,E').
\]
These functors are compatible, up to their canonical associativity
isomorphisms, with successive finite extensions of coefficient fields.
\end{theorem}

\begin{definition}[Ambient pullback, tensor product, and direct image]
For a morphism \(f:X\to Y\), the derived operations used below are
\[
 Lf^*:D(Y,E)\longrightarrow D(X,E),\qquad
 Rf_*:D(X,E)\longrightarrow D(Y,E),
\]
and
\[
 -\otimes_E^L-:D(X,E)\times D(X,E)\longrightarrow D(X,E).
\]
Derived internal Hom is the contravariant triangulated functor
\[
 R\mathcal Hom_X:
 D(X,E)^{\mathrm{op}}\times D(X,E)\longrightarrow D(X,E).
\]
They are computed using \(K\)-flat and \(K\)-injective resolutions.  Their
restrictions to bounded constructible categories are supplied by the
corresponding finiteness theorems.
\end{definition}

\begin{definition}[Compactly supported direct image]
Let \(f:X\to Y\) be separated and of finite type.  Choose a compactification
\[
             X\mathop{\lhook\joinrel\longrightarrow}^{j}\overline X
              \mathop{\longrightarrow}^{\overline f}Y
\]
with \(j\) an open immersion and \(\overline f\) proper.  Extension by zero
is the exact functor \(j_!\) on sheaves, and
\[
             Rf_!:=R\overline f_*\,j_!:D(X,E)\longrightarrow D(Y,E).
\]
The compactification-independence theorem supplies canonical comparisons
for different choices; these comparisons satisfy identity, transitivity,
and composition compatibility.
\end{definition}

\begin{theorem}[Finiteness for the operations used below]
For a finite-type morphism \(f:X\to Y\), the functors \(Lf^*\) and \(Rf_*\)
send bounded constructible complexes to bounded constructible complexes.  If
\(f\) is separated, the same holds for \(Rf_!\).  Derived tensor product and
derived internal Hom preserve bounded constructibility on \(X\).  If
\(j:U\hookrightarrow C\) is a dense open
immersion of smooth curves and \(\mathcal F\) is lisse of finite rank on
\(U\), then \(j_*\mathcal F=R^0j_*\mathcal F\) and \(R^1j_*\mathcal F\)
are constructible, \(R^ij_*\mathcal F=0\) for \(i>1\), and
\(Rj_*\mathcal F\in\Dbc(C,E)\).  Thus all middle extensions and Fourier
transforms below lie in the displayed bounded constructible categories.
\end{theorem}

\begin{theorem}[Scalar extension and the derived operations]
Let \(E'/E\) be finite, let \(f:X\to Y\) be a finite-type morphism, and
take \(A\in\Dbc(Y,E)\) and \(K,L\in\Dbc(X,E)\).  There are canonical
isomorphisms
\[
\begin{aligned}
 (Lf^*A)\otimes_E^LE'
   &\simeq Lf^*(A\otimes_E^LE'),\\
 (Rf_*K)\otimes_E^LE'
   &\simeq Rf_*(K\otimes_E^LE'),\\
 (K\otimes_E^LL)\otimes_E^LE'
   &\simeq (K\otimes_E^LE')\otimes_{E'}^L(L\otimes_E^LE'),\\
 R\mathcal Hom_X(K,L)\otimes_E^LE'
   &\simeq R\mathcal Hom_X(K\otimes_E^LE',L\otimes_E^LE').
\end{aligned}
\]
If \(f\) is separated, there is in addition a canonical isomorphism
\[
 (Rf_!K)\otimes_E^LE'\simeq Rf_!(K\otimes_E^LE').
\]
These isomorphisms are functorial in the displayed arguments, are compatible
with successive coefficient extensions and composition of morphisms, and in
the compact-support case respect the compactification-independence
comparisons.  Consequently the operations just defined act on
\(\mathcal C_X\) and on its Karoubi completion \(\Dbc(X,\Qlb)\).
\end{theorem}

\begin{definition}[Compactly supported cohomology]
For the structure morphism \(a_X:X\to\operatorname{Spec}\Fbar\), define
\[
 R\Gamma_c(X,K)=Ra_{X!}K,\qquad
 H_c^i(X,K)=\mathcal H^i(Ra_{X!}K).
\]
If \(K\) has a Weil structure, compact-support base change transports it to
these finite-dimensional cohomology groups.
\end{definition}

\section{Weil sheaves and weights}

Write \(X=X_0\times_k\Fbar\).  Let
\[
 \sigma_{\mathrm{ar}}(a)=a^q,\qquad
 \sigma_{\mathrm{geom}}=\sigma_{\mathrm{ar}}^{-1}.
\]
Let \(\gamma_X:X\to X\) be the semilinear automorphism induced by
\(\sigma_{\mathrm{geom}}\).  The \(q\)-power scheme Frobenius
\(\operatorname{Fr}_{X_0,q}:X_0\to X_0\), locally \(f\mapsto f^q\), is a
different morphism and is not denoted by \(\gamma_X\).

\begin{definition}[Weil sheaf]
A Weil sheaf on \(X_0\) is a constructible sheaf \(\mathcal F\) on \(X\)
together with an isomorphism
\[
 \varphi:\gamma_X^*\mathcal F\xrightarrow{\sim}\mathcal F.
\]
Iterating this datum at a closed point defines the operator denoted
\(\Frob_x\) on \(\mathcal F_{\bar x}\); it is geometric Frobenius
\(\sigma_{\mathrm{geom}}^{\deg(x)}\), not arithmetic Frobenius.  The
 operators induced by the six operations, in particular on compactly
supported cohomology, use the same convention and do not arise by applying
cohomology to \(\operatorname{Fr}_{X_0,q}\).  A lisse Weil sheaf is one
whose underlying geometric sheaf is locally constant of finite rank.  Write
\(\operatorname{LisseW}(X_0,E)\) for the category of lisse Weil sheaves;
its morphisms are morphisms of the underlying lisse sheaves that commute
with the displayed Weil isomorphisms.
\end{definition}

\begin{definition}[Weil complexes by equivariant localization]
Use the canonical identification \(\gamma_X^{-1}E_X\simeq E_X\) when
forming pullback of coefficient modules.
Let \(\mathsf C^b_{c,W}(X_0,E)\) be the category whose objects are bounded
complexes \(K^\bullet\) of \(E_X\)-module sheaves with constructible
cohomology, equipped with an isomorphism of complexes
\[
       \phi_K:\gamma_X^*K^\bullet\xrightarrow{\sim}K^\bullet.
\]
A morphism is a chain map commuting with the displayed isomorphisms.  An
equivariant chain homotopy is a chain homotopy commuting with the same
isomorphisms.  Let \(\mathsf K^b_{c,W}(X_0,E)\) be the resulting homotopy
category.  Shifts and mapping cones inherit strict Weil data and make this
an ordinary triangulated category.  An equivariant quasi-isomorphism is a
morphism whose underlying chain map is a quasi-isomorphism.  Define the
bounded constructible Weil derived category by the Verdier localization
\[
 \DbcW(X_0,E)=
 \mathsf K^b_{c,W}(X_0,E)
 [\{\text{equivariant quasi-isomorphisms}\}^{-1}].
\]
Its distinguished triangles are those induced from the equivariant homotopy
category.  Forgetting the strict datum before localization gives an exact functor
\[
       \operatorname{For}_X:\DbcW(X_0,E)\longrightarrow\Dbc(X,E).
\]
A constructible Weil sheaf is regarded as an object of this category
concentrated in degree zero.
\end{definition}

\begin{theorem}[Lift of geometric operations to Weil complexes]
Pullback, tensor product, ordinary and compactly supported direct image, and
the Fourier transforms used below admit exact lifts to the corresponding
categories \(\DbcW\).  After applying \(\operatorname{For}\), the lifted
functor is canonically isomorphic to the underlying geometric derived
operation.  Its Weil datum is induced by the base-change comparison for
\(\gamma\), together with the Weil datum on the Artin--Schreier kernel in
the Fourier case.  The lifts commute with finite scalar extension and with
composition wherever the underlying derived operations do.
\end{theorem}

\begin{definition}[Strict Weil complexes with \(\Qlb\)-coefficients]
Let \(\mathcal C^W_{X_0}\) have objects \((E,K)\), where \(E/\Ql\) is
finite and \(K\in\DbcW(X_0,E)\).  For two objects define
\[
 \operatorname{Hom}_{\mathcal C^W_{X_0}}((E,K),(F,L))
 =\varinjlim_{M\supset E,F}
 \operatorname{Hom}_{\DbcW(X_0,M)}
 \bigl(K\otimes_E^LM,L\otimes_F^LM\bigr),
\]
where \(M\) runs through common finite extensions of \(\Ql\) in \(\Qlb\).
Composition is performed after scalar extension to one common stage.  Put
\[
                    \DbcW(X_0,\Qlb)=(\mathcal C^W_{X_0})^\natural.
\]
Thus an object is a triple \((E,K,e)\) with \(e^2=e\), a morphism
\((E,K,e)\to(F,L,e')\) is a common-stage equivariant class \(u\) satisfying
\(u=e'u=ue\), and distinguished triangles are direct summands of
equivariant distinguished triangles represented over a common finite
extension of \(\Ql\).
\end{definition}

\begin{theorem}[Operations on \(\Qlb\)-Weil complexes]
The lifted geometric operations commute with finite scalar extension and
therefore induce exact triangulated functors on \(\DbcW(-,\Qlb)\).  For a
Fourier transform with fixed character \(\psi\), the common stages are
restricted to coefficient extensions containing the values of \(\psi\).
The forgetful functors to \(\Dbc(-,\Qlb)\) commute with all these operations.
\end{theorem}

\begin{definition}[Tate twists]
On the pro-etale site put
\[
 \mathbf Z_\ell(1)_X=\varprojlim_n\mu_{\ell^n,X},
 \qquad
 E_X(1)=
 \left(\mathbf Z_\ell(1)_X
 \otimes_{(\mathbf Z_\ell)_X}(\mathcal O_E)_X\right)[1/\varpi].
\]
For \(m\ge0\), let \(E_X(m)=E_X(1)^{\otimes m}\), and for \(m<0\) use
the corresponding tensor power of the dual line \(E_X(1)^\vee\).  For
\(K\in D(X,E)\), define
\[
                         K(m)=K\otimes^L_{E_X}E_X(m).
\]
The canonical descent datum on \(E_X(m)\) defines the same operation on
strict Weil complexes.  At a closed point \(x\), geometric Frobenius acts
on \(E_X(1)_{\bar x}\) by \(q_x^{-1}\); consequently \((m)\) lowers weights
by \(2m\).

For a local decomposition or Weil group \(G\) with residue cardinality
\(Q\), let \(E(1)\) be the one-dimensional cyclotomic representation on
which geometric Frobenius acts by \(Q^{-1}\).  Put
\(V(m)=V\otimes_EE(1)^{\otimes m}\), using the dual line when \(m<0\).
This agrees with the stalk of the sheaf-theoretic twist.
\end{definition}

\begin{definition}[Geometric Frobenius and power trace functions]
Let \(X_0\) be a finite-type \(k\)-scheme and let
\(K_0\in\DbcW(X_0,E)\).  Write \(K=\operatorname{For}_X(K_0)\) and let
\(\phi_K:\gamma_X^*K\xrightarrow{\sim}K\) be its strict Weil isomorphism.
For \(n\geq1\), define
\[
 \phi_K^{(n)}=
 \phi_K\circ\gamma_X^*(\phi_K)\circ\cdots\circ
 \gamma_X^{(n-1)*}(\phi_K):\gamma_X^{n*}K\xrightarrow{\sim}K.
\]
If \(x\in X_0(\mathbf F_{q^n})\), choose a geometric point \(\bar x\)
above \(x\).  The \(\mathbf F_{q^n}\)-rational structure identifies
\(\gamma_X^n(\bar x)\) with \(\bar x\).  On stalk cohomology the map
\(\phi_K^{(n)}\) induces
\[
 F_{x,n}^i(K_0)\in
 \operatorname{Aut}_E\bigl(\mathcal H^i(K)_{\bar x}\bigr).
\]
If \(|x|\) is the closed point below \(x\), of degree \(d\mid n\), then,
up to conjugacy from the geometric lift,
\[
 F_{x,n}^i(K_0)=\bigl(F_{|x|}^i(K_0)\bigr)^{n/d}
       \quad\text{on }\mathcal H^i(K)_{\bar x}.
\]
The degree-\(n\) trace function is
\[
 \begin{aligned}
 t_{K_0}^{(n)}:X_0(\mathbf F_{q^n})&\longrightarrow E,\\
 x&\longmapsto
 \sum_i(-1)^i\operatorname{Tr}\!\left(
 F_{x,n}^i(K_0)\mid\mathcal H^i(K)_{\bar x}\right).
 \end{aligned}
\]
It is independent of the chosen geometric lift.  For \(\Qlb\)-coefficients
the operator and trace are represented at one common finite coefficient
stage and are unchanged by further scalar extension.
\end{definition}

\begin{definition}[Arithmetic and geometric irreducibility]
A lisse Weil sheaf is \emph{arithmetically irreducible} if it has no nonzero
proper subsheaf stable under its Weil structure.  It is
\emph{geometrically irreducible} if its underlying lisse sheaf on
\(X_0\times_k\Fbar\) is irreducible.  Geometric irreducibility implies
arithmetic irreducibility.  The converse is not assumed.
\end{definition}

\begin{definition}[Purity and weight inequalities]
Put \(q_x=|k(x)|\).  A lisse Weil sheaf \(\mathcal F\) is punctually
\(\iota\)-pure of weight \(w\in\mathbf R\) if every eigenvalue \(\alpha\)
of \(\Frob_x\) on \(\mathcal F_{\bar x}\), for every closed point \(x\),
satisfies
\[
                         |\iota(\alpha)|=q_x^{w/2}.
\]
For a finite-dimensional vector space \(M\), an automorphism \(F\), and a
specified real number \(Q>1\), the Frobenius module \((M,F;Q)\) has weights
at most \(v\) if
\[
 |\iota(\alpha)|\le Q^{v/2}
\]
for every eigenvalue \(\alpha\) of \(F\), and it has weights at least \(v\)
if all these inequalities are reversed.  It is pure of weight \(v\) if
equality holds for every eigenvalue.
The Tate twist \(V(m)\) is normalized to lower weights by \(2m\); in
particular, \(V(-1)\) has weight two more than \(V\).
\end{definition}

\begin{definition}[Real and conjugate Weil sheaves]
A lisse Weil sheaf \(\mathcal G\) is \emph{\(\iota\)-real} if
\[
 \iota\!\left(\det(1-T\Frob_x\mid\mathcal G_{\bar x})\right)
                         \in\mathbf R[T]
\]
for every closed point \(x\).

An \emph{\(\iota\)-conjugate companion} of \(\mathcal F\) is a lisse Weil
sheaf \(\mathcal F^c\) such that, at every closed point, the Frobenius
eigenvalue multiset of \(\mathcal F^c\), after applying \(\iota\), is the
complex-conjugate multiset of that of \(\mathcal F\).  This definition is
local at every closed point; equality of a single global Euler product is
not sufficient.
\end{definition}

\section{Local monodromy}

Let \(C_0\) be a smooth curve over \(k\), let \(s\) be a closed point, and
put
\[
 T_s=\operatorname{Spec}(\mathcal O^h_{C_0,s}),\qquad
 \eta_s=\operatorname{Spec}(\operatorname{Frac}\mathcal O^h_{C_0,s}).
\]
Choose a geometric generic point \(\bar\eta_s\).

\begin{definition}[Decomposition and inertia groups]
The local decomposition and inertia groups are
\[
 D_s=\pi_1(\eta_s,\bar\eta_s),\qquad
 I_s=\ker\!\left(D_s\longrightarrow
          \operatorname{Gal}(\overline{k(s)}/k(s))\right).
\]
They occur in the canonical exact sequence
\[
 1\longrightarrow I_s\longrightarrow D_s\longrightarrow
 \operatorname{Gal}(\overline{k(s)}/k(s))\longrightarrow1.
\]
Let \(W_s\subset D_s\) be the inverse image of the infinite cyclic subgroup
generated by geometric Frobenius in the residue-field Galois group.  It is
the local Weil group and fits into
\[
 1\longrightarrow I_s\longrightarrow W_s\longrightarrow
 \langle\Frob_s\rangle\simeq\mathbf Z\longrightarrow1.
\]
The wild inertia group \(P_s\) is the unique pro-\(p\) Sylow subgroup of
\(I_s\), and \(I_s^{(r)}\), \(r\ge0\), denotes the upper-numbering
ramification filtration.
\end{definition}

\begin{definition}[Local representation category]
The category \(\Rep_E(D_s)\) has finite-dimensional \(E\)-vector spaces with
continuous \(D_s\)-action as objects and equivariant linear maps as
morphisms.  Write \(\operatorname{Lisse}(\eta_s,E)\) for the category of
lisse finite-rank \(E\)-sheaves on \(\eta_s\).

The category \(\Rep_E(W_s)\) consists of finite-dimensional
\(E\)-representations whose restriction to the open subgroup \(I_s\) is
continuous; \(W_s/I_s\simeq\mathbf Z\) is given the discrete topology.
Write \(\operatorname{LisseW}(\eta_s,E)\) for the category of lisse
finite-rank sheaves on the geometric generic point equipped with a strict
local Weil descent datum.  Restriction of the group action defines an exact
functor
\[
 \operatorname{Res}^{D_s}_{W_s}:\Rep_E(D_s)
                  \longrightarrow\Rep_E(W_s).
\]
\end{definition}

\begin{theorem}[Local monodromy equivalences]
Geometric generic stalks define exact equivalences of abelian categories
\[
 \begin{aligned}
 (-)_{\bar\eta_s}:\operatorname{Lisse}(\eta_s,E)
     &\xrightarrow{\sim}\Rep_E(D_s),\\
 (-)_{\bar\eta_s}:\operatorname{LisseW}(\eta_s,E)
     &\xrightarrow{\sim}\Rep_E(W_s).
 \end{aligned}
\]
Denote chosen quasi-inverses by
\[
 \operatorname{Loc}_s^D:\Rep_E(D_s)\longrightarrow
      \operatorname{Lisse}(\eta_s,E),\qquad
 \operatorname{Loc}_s^W:\Rep_E(W_s)\longrightarrow
      \operatorname{LisseW}(\eta_s,E).
\]
If \(\operatorname{For}^{D_s}_{W_s}\) sends an arithmetic lisse sheaf to
its strict local Weil sheaf, then there is a functorial isomorphism
\[
 \operatorname{For}^{D_s}_{W_s}\!\left(
       \operatorname{Loc}_s^D(V)\right)
 \xrightarrow{\sim}
 \operatorname{Loc}_s^W\!\left(
       \operatorname{Res}^{D_s}_{W_s}V\right)
 \qquad(V\in\Rep_E(D_s)).
\]
The two equivalences, their unit and counit isomorphisms, and this
comparison commute with finite extension of \(E\).
\end{theorem}

\begin{definition}[Breaks and Swan conductor]
For \(V\in\Rep_E(W_s)\), the restriction to wild inertia has a canonical
break decomposition
\[
                         V=\bigoplus_{\lambda\ge0}V(\lambda).
\]
The finitely many \(\lambda\) for which \(V(\lambda)\ne0\) are the breaks,
or slopes, and
\[
                 \Sw(V)=\sum_{\lambda\ge0}\lambda\dim_E V(\lambda).
\]
For an interval \(J\subset\mathbf R_{\ge0}\), let
\(\Rep_E(W_s)_J\) be the full abelian subcategory whose objects have all
breaks in \(J\).  Put
\[
 \Rep_E(D_s)_J=\{V\in\Rep_E(D_s):
       \operatorname{Res}^{D_s}_{W_s}V\in\Rep_E(W_s)_J\}.
\]
A representation of either local group is tame when all breaks are zero and
is unramified when \(I_s\) acts trivially.
\end{definition}

\begin{theorem}[Break decomposition and integrality of Swan conductor]
For every \(V\in\Rep_E(W_s)\), there is a unique finite direct-sum
decomposition
\[
 V=\bigoplus_{\lambda\in\mathbf Q_{\ge0}}V(\lambda)
\]
such that, for every \(r>0\),
\[
 V(\lambda)^{I_s^{(r)}}=
 \begin{cases}
 0,&0<r\le\lambda,\\
 V(\lambda),&r>\lambda.
 \end{cases}
\]
The decomposition is functorial in \(V\), exact on short exact sequences,
and preserved by finite extension of \(E\).  Every break is rational and
\[
 \Sw(V)=\sum_\lambda\lambda\dim_E V(\lambda)\in\mathbf Z_{\ge0}.
\]
For a direct sum, Swan conductors add; for an unramified representation the
Swan conductor is zero.  If \(V\in\Rep_E(D_s)\), this decomposition and
Swan conductor agree with those of
\(\operatorname{Res}^{D_s}_{W_s}V\).
\end{theorem}

\begin{definition}[Invariants and echelon-two monodromy]
For \(V\in\Rep_E(W_s)\), define
\[
 V^{I_s}=\{v\in V:\sigma v=v\text{ for every }\sigma\in I_s\},
 \qquad W_1(V)=V/V^{I_s}.
\]
The inertia action is \emph{unipotent of echelon two} if \(I_s\) acts
trivially on \(W_1(V)\).  Equivalently,
\((\sigma-1)V\subset V^{I_s}\) for every \(\sigma\in I_s\), and hence
\((\sigma-1)^2=0\).  In this case both ends of
\[
                  0\longrightarrow V^{I_s}\longrightarrow V
                    \longrightarrow W_1(V)\longrightarrow0
\]
are representations of the residue-field Galois group through
\(W_s/I_s\).
\end{definition}

\section{Middle extension and type 3}

Let \(A_0=\Aone_k\), let \(A=A_0\times_k\Fbar\), and let
\(j_0:U_0\hookrightarrow A_0\) be dense, with geometric base change
\(j:U\hookrightarrow A\).

\begin{definition}[Artin--Schreier sheaf]
Let
\[
             \operatorname{AS}:\Aone_k\longrightarrow\Aone_k,
             \qquad z\longmapsto z^q-z .
\]
This is a finite etale Galois covering with Galois group the additive
group of \(k\), acting by translations on the source.  The
\(\psi\)-isotypic rank-one summand of
\(\operatorname{AS}_*E_{\Aone_k}\) is denoted
\(\mathcal L_\psi^{\mathrm{arith}}\).  Its pullback to \(A\), together
with the induced strict Weil datum, is denoted \(\mathcal L_\psi\).  For
every closed point \(x\) of \(A_0\),
\[
 \operatorname{Tr}\!\left(
   \Frob_x\mid(\mathcal L_\psi)_{\bar x}\right)
   =\psi\!\left(\operatorname{Tr}_{k(x)/k}(x)\right).
\]
Both sheaves are lisse of rank one and punctually pure of weight zero, and
there are canonical duality isomorphisms
\[
 (\mathcal L_\psi^{\mathrm{arith}})^\vee
     \simeq\mathcal L_{\psi^{-1}}^{\mathrm{arith}},
 \qquad
 \mathcal L_\psi^\vee\simeq\mathcal L_{\psi^{-1}}.
\]
\end{definition}

\begin{definition}[Verdier duality on the affine line]
Let \(E_A\) be the completed rational pro-etale coefficient sheaf.  The
dualizing complex on the smooth curve \(A\) is
\(\omega_A=E_A(1)[2]\).  Define
\[
 \begin{aligned}
 D_A:\Dbc(A,E)^{\mathrm{op}}&\longrightarrow\Dbc(A,E),\\
 K&\longmapsto R\mathcal Hom_A(K,\omega_A).
 \end{aligned}
\]
\end{definition}

\begin{theorem}[Constructible Verdier duality]
For every \(K\in\Dbc(A,E)\), the ambient internal-Hom complex
\(R\mathcal Hom_A(K,\omega_A)\) has bounded constructible cohomology.
The functor \(D_A\) is an exact contravariant equivalence, and the
evaluation morphism \(K\to D_AD_AK\) is an isomorphism.  This statement
commutes with finite coefficient extension and descends to \(\Qlb\).
\end{theorem}

\begin{definition}[Geometric perverse middle extension]
The middle perverse \(t\)-structure is given by
\[
\begin{aligned}
 {}^pD^{\le0}(A,E)
 &=\{K\in\Dbc(A,E):\dim\operatorname{Supp}\mathcal H^m(K)\le -m
       \text{ for every }m\},\\
 {}^pD^{\ge0}(A,E)
 &=\{K\in\Dbc(A,E):D_AK\in{}^pD^{\le0}(A,E)\}.
\end{aligned}
\]
Its heart is
\(\Perv(A,E)={}^pD^{\le0}(A,E)\cap{}^pD^{\ge0}(A,E)\), and
\({}^pH^0\) is degree-zero perverse cohomology.  If \(\mathcal F\) is
lisse on \(U\), then \(\mathcal F[1]\) is perverse and
\[
 j_{!*}(\mathcal F[1])=
 \operatorname{im}_{\Perv(A,E)}\!\left(
 {}^pH^0(j_!\mathcal F[1])\longrightarrow
 {}^pH^0(Rj_*\mathcal F[1])\right).
\]
On a smooth curve there is a canonical isomorphism
\[
                       j_{!*}(\mathcal F[1])\simeq(j_*\mathcal F)[1],
\]
where \(j_*\mathcal F=R^0j_*\mathcal F\) is the ordinary direct image.
This is not an identification with \(Rj_*\mathcal F[1]\).
\end{definition}

\begin{definition}[Strict Weil perverse truncation subcategories]
Put
\[
 \begin{aligned}
 {}^pD_W^{\leq0}(A_0,E)
 &=\{K\in\DbcW(A_0,E):\operatorname{For}_A(K)\in{}^pD^{\leq0}(A,E)\},\\
 {}^pD_W^{\geq0}(A_0,E)
 &=\{K\in\DbcW(A_0,E):\operatorname{For}_A(K)\in{}^pD^{\geq0}(A,E)\}.
 \end{aligned}
\]
\end{definition}

\begin{theorem}[Strict Weil perverse \(t\)-structure]
The pair
\[
 \bigl({}^pD_W^{\leq0}(A_0,E),{}^pD_W^{\geq0}(A_0,E)\bigr)
\]
is a bounded \(t\)-structure on \(\DbcW(A_0,E)\).  Its heart
\[
 \PervW(A_0,E)=
 {}^pD_W^{\leq0}(A_0,E)\cap{}^pD_W^{\geq0}(A_0,E)
\]
is an abelian category.  Write \({}^pH_W^n\) for its cohomology functors.
The forgetful functor is perverse \(t\)-exact, and for every
\(K\in\DbcW(A_0,E)\) there are functorial isomorphisms
\[
 \operatorname{For}_A({}^pH_W^nK)
       \simeq{}^pH^n(\operatorname{For}_A K).
\]
\end{theorem}

\begin{theorem}[Scalar extension on strict Weil perverse hearts]
For every finite extension \(E'/E\), the exact functor
\[
 -\otimes_E^LE':\DbcW(A_0,E)\longrightarrow\DbcW(A_0,E')
\]
is perverse \(t\)-exact.  It therefore induces an exact functor
\[
 \PervW(A_0,E)\longrightarrow\PervW(A_0,E').
\]
It commutes with \({}^pH_W^n\), with the forgetful functor, and with the
maps between perverse cohomology objects induced by \(j_{0!}\to Rj_{0*}\).
Because \(E'/E\) is flat, it also commutes with kernels, cokernels, and
images in the two perverse hearts.
\end{theorem}

\begin{definition}[Strict Weil perverse middle extension]
Let \(j_0:U_0\hookrightarrow A_0\) be a dense open immersion, and let
\(\mathcal F_0\) be a lisse Weil \(E\)-sheaf on \(U_0\).  Define
\[
 j_{0!*}^W(\mathcal F_0[1])=
 \operatorname{im}_{\PervW(A_0,E)}\!\left(
 {}^pH_W^0(j_{0!}\mathcal F_0[1])\longrightarrow
 {}^pH_W^0(Rj_{0*}\mathcal F_0[1])\right).
\]
\end{definition}

\begin{theorem}[Forgetting the strict Weil middle extension]
With the notation of the preceding definition, put
\(\mathcal F=\operatorname{For}_{U}(\mathcal F_0)\).  There is a canonical
isomorphism in \(\Perv(A,E)\)
\[
 \operatorname{For}_A\bigl(j_{0!*}^W(\mathcal F_0[1])\bigr)
       \simeq j_{!*}(\mathcal F[1]).
\]
This comparison is functorial in \(\mathcal F_0\) and commutes with finite
extension of the coefficient field.
\end{theorem}

\begin{theorem}[Common-stage and Karoubi perverse hearts]
Let \(\mathcal P_A\) be the category whose objects are pairs \((E,P)\) with
\(P\in\Perv(A,E)\) and whose morphisms are
\[
 \operatorname{Hom}_{\mathcal P_A}((E,P),(F,Q))=
 \varinjlim_{M\supset E,F}
 \operatorname{Hom}_{\Perv(A,M)}
 (P\otimes_E M,Q\otimes_F M).
\]
Define \(\mathcal P^W_{A_0}\) by the same formula with
\(\PervW(A_0,-)\) in place of \(\Perv(A,-)\).  In
\(\Dbc(A,\Qlb)\), respectively \(\DbcW(A_0,\Qlb)\), let the nonpositive
and nonnegative subcategories consist of objects isomorphic to direct
summands of finite-stage objects lying in the corresponding finite-stage
perverse truncation subcategory.  These pairs are bounded
\(t\)-structures, and their abelian hearts are canonically equivalent to
\[
 \Perv(A,\Qlb)\simeq\mathcal P_A^\natural,
 \qquad
 \PervW(A_0,\Qlb)\simeq(\mathcal P^W_{A_0})^\natural.
\]
The finite-stage functors into these hearts are exact.  The forgetful
functor between the two \(\Qlb\)-coefficient derived categories is
perverse \(t\)-exact.  The finite-stage middle extensions are compatible
with scalar extension and therefore determine middle-extension objects in
these two Karoubi-completed hearts; the forgetful comparison of the
preceding theorem remains valid over \(\Qlb\).
\end{theorem}

\begin{definition}[Geometric and arithmetic type 3]
For \(a\in\Fbar\), let \(m_a:A\to A\) be multiplication by \(a\), and put
\(\mathcal L_{\psi,a}=m_a^*\mathcal L_\psi\); the case \(a=0\) is the
constant rank-one sheaf.  A simple object of \(\Perv(A,\Qlb)\) is of:
\begin{enumerate}
\item type 1 if it is supported at a closed point;
\item type 2 if it is \(\mathcal L_{\psi,a}[1]\) for some
\(a\in\Fbar\);
\item type 3 if it is \(j_{!*}(\mathcal G[1])\) for an irreducible lisse
\(\Qlb\)-sheaf \(\mathcal G\) on a dense open and \(\mathcal G\) is not
isomorphic, after restriction to any smaller dense open, to any
\(\mathcal L_{\psi,a}\).
\end{enumerate}
For an arithmetically irreducible lisse Weil \(E\)-sheaf \(\mathcal F_0\),
put \(\mathcal K_{\mathcal F_0}=j_{0!*}^W(\mathcal F_0[1])\).  It is
\emph{arithmetically of type 3} if every simple constituent of
\[
 \operatorname{For}_A(\mathcal K_{\mathcal F_0})
       \otimes_E^L\Qlb\in\Perv(A,\Qlb)
\]
is geometrically of type 3.  Arithmetic type 3 does not assert geometric
irreducibility.
\end{definition}

\begin{lemma}[Artin--Schreier alternatives at infinity]
The sheaf \(\mathcal L_{\psi,0}\) is constant.  For every
\(a\in\Fbar^\times\), the sheaf \(\mathcal L_{\psi,a}\) has Swan conductor
one at infinity, and its restriction to every dense open of \(A\) remains
ramified at infinity.
\end{lemma}

\begin{proposition}[A useful type-3 criterion]
Let \(\mathcal F_0\) be an arithmetically irreducible lisse Weil sheaf on
\(U_0\).  Suppose its geometric local monodromy is unramified at infinity
and its underlying geometric local system has no constant simple
subquotient.  No geometric simple constituent of \(\mathcal F_0\) is
isomorphic, on any smaller dense open, to \(\mathcal L_{\psi,a}\) for any
\(a\in\Fbar\).  Every simple constituent of
\[
 \operatorname{For}_A(\mathcal K_{\mathcal F_0})\otimes_E^L\Qlb
\]
is geometrically of type 3, and \(\mathcal K_{\mathcal F_0}\) is
arithmetically of type 3.
\end{proposition}

\section{The global Fourier--Deligne transform}
\label{sec:global-ft}

Let \(A'_0\) be the line dual to \(A_0\), let \(A'\) be its geometric base
change, and write
\[
 p:A\times A'\to A,\qquad p':A\times A'\to A',\qquad \mu(x,x')=xx'.
\]
Write \(p_0,p'_0:A_0\times_kA'_0\to A_0,A'_0\) for the arithmetic
projections and \(\mu_0:A_0\times_kA'_0\to\Aone_k\) for multiplication.

\begin{definition}[Fourier--Deligne transform]
The normalized Fourier--Deligne transform is the triangulated functor
\[
\begin{aligned}
 \FT_\psi:\Dbc(A,E)&\longrightarrow\Dbc(A',E),\\
 K&\longmapsto
 Rp'_!\bigl(Lp^*K\otimes_E^L L\mu^*\mathcal L_\psi\bigr)[1].
\end{aligned}
\]
The shift \([1]\) is part of the definition, and there is no defining Tate
twist.  The same formula on the arithmetic lines defines
\[
 \FT_{\psi,0}:\Dbc(A_0,E)\longrightarrow\Dbc(A'_0,E),
 \qquad
 K_0\longmapsto
 Rp'_{0!}\bigl(Lp_0^*K_0\otimes_E^L
 L\mu_0^*\mathcal L_\psi^{\mathrm{arith}}\bigr)[1].
\]
It commutes with geometric base change and finite extension of coefficients.
The lifted functor on Weil complexes is denoted
\[
 \FT_{\psi,W}:\DbcW(A_0,E)\longrightarrow\DbcW(A'_0,E).
\]
\end{definition}

\begin{theorem}[Stalk formula]
For \(K\in\Dbc(A,E)\) and a geometric point \(\bar y\to A'\), let
\(\mathcal L_{\psi,y}\) be the pullback of \(\mathcal L_\psi\) along
\(x\mapsto xy\).  There is a functorial isomorphism
\[
 (\FT_\psi K)_{\bar y}\simeq
       \RGammac(A,K\otimes_E^L\mathcal L_{\psi,y})[1].
\]
If \(K\) is the geometric pullback of a Weil complex on \(A_0\) and
\(\bar y\) lies over a closed point \(y_0\in A'_0\), this comparison is
equivariant for the geometric Frobenius at \(y_0\).
\end{theorem}

\begin{theorem}[Perverse equivalence and type exchange]
The functor \(\FT_\psi\) is perverse \(t\)-exact and induces an exact
equivalence
\[
                     \Perv(A,E)\xrightarrow{\sim}\Perv(A',E).
\]
After finite scalar extension and Karoubi completion, it exchanges types 1
and 2 in \(\Perv(A,\Qlb)\) and preserves geometric type 3.  The strict
Weil transform \(\FT_{\psi,W}\) is perverse \(t\)-exact and induces an
exact equivalence
\[
             \PervW(A_0,E)\xrightarrow{\sim}\PervW(A'_0,E),
\]
compatible with the forgetful functor.  It carries arithmetic type-3
objects to arithmetic type-3 objects.  In particular, the transform of a
simple geometric type-3 middle extension is a whole perverse object
\(v_{!*}(\mathcal H[1])\), where \(v:V\hookrightarrow A'\) is a dense
open immersion and \(\mathcal H\) is irreducible and lisse on \(V\); it
has no punctual direct summand.
\end{theorem}

\begin{theorem}[Laumon's universal support-forgetting comparison]
For every \(K\in D^b_c(A,E)\), the canonical morphism
\[
 Rp'_!\bigl(Lp^*K\otimes_E^L L\mu^*\mathcal L_\psi\bigr)
 \longrightarrow
 Rp'_*\bigl(Lp^*K\otimes_E^L L\mu^*\mathcal L_\psi\bigr)
\]
is an isomorphism in \(D^b_c(A',E)\).  It is compatible with finite
extension of the coefficient field.  Consequently \(\FT_\psi\) may
equivalently be written
using \(Rp'_*\), although its definition remains the compactly supported
formula above.
\end{theorem}

\section{The three local Fourier transforms}

Let
\[
 R=(k[\pi]_{(\pi)})^h,\qquad
 R'=(k[\pi']_{(\pi')})^h
\]
over the fixed finite field \(k\), and put
\(T=\operatorname{Spec}R\) and \(T'=\operatorname{Spec}R'\).  Their closed
points are denoted by \(t,t'\), their generic points by \(\eta,\eta'\), and
their decomposition groups by \(D_T,D_{T'}\).  Write
\(\operatorname{pr},\operatorname{pr}'\) for the projections from
\(T\times_kT'\).  The chosen coordinates define involutions
\[
 a_T:T\longrightarrow T,\quad \pi\longmapsto-\pi,
 \qquad
 a_{T'}:T'\longrightarrow T',\quad \pi'\longmapsto-\pi'.
\]

For a rational function \(f\), write
\(\overline{\mathcal L}_\psi(f)\) for extension by zero to
\(T\times_kT'\) of the corresponding Artin--Schreier sheaf.  For
\(f=\pi/\pi'\), \(f=\pi'/\pi\), and \(f=1/(\pi\pi')\), its respective
domains are \(T\times_k\eta'\), \(\eta\times_kT'\), and
\(\eta\times_k\eta'\).

\begin{definition}[Equivariant cycle categories]
Put \(Z=T\times_kT'\) and \(g=\operatorname{pr}':Z\to T'\).  Choose a
geometric generic point \(\bar\eta'\), a geometric closed point
\(\bar t'\), and a specialization \(\bar\eta'\to\bar t'\).  If
\(T'^{\mathrm{sh}}\) is the strict henselization at \(\bar t'\), put
\(Z^{\mathrm{sh}}=Z\times_{T'}T'^{\mathrm{sh}}\), and write
\[
 \bar j:Z_{\bar\eta'}\longrightarrow Z^{\mathrm{sh}},\qquad
 \bar i:Z_{\bar t'}\longrightarrow Z^{\mathrm{sh}}.
\]
For \(Y=Z_{\bar\eta'}\) or \(Z_{\bar t'}\), let
\(\Dbc(Y,E;D_{T'})\) be the ordinary Verdier localization of bounded
constructible complexes carrying continuous \(D_{T'}\)-equivariant
descent data for the natural action on \(Y\).  Define
\(\Dbc(Y,E;W_{T'})\) in the same way with \(W_{T'}\): continuity is
required on inertia, and a geometric Frobenius operator satisfies the Weil
conjugation relation.  Restriction of equivariance gives exact
triangulated functors
\[
 \operatorname{Res}^{D_{T'}}_{W_{T'}}:
 \Dbc(Y,E;D_{T'})\longrightarrow\Dbc(Y,E;W_{T'}).
\]
Finally, \(\DbcW(Z,E)\) is the ordinary equivariant Verdier localization
of bounded constructible complexes on \(Z_{\Fbar}\) with a strict
geometric-Frobenius isomorphism.  Pullback of an arithmetic complex and
restriction of its descent datum define
\[
 \operatorname{For}^{D}_{W,Z}:\Dbc(Z,E)\longrightarrow\DbcW(Z,E).
\]
\end{definition}

\begin{definition}[Nearby and vanishing cycles with typed actions]
For \(H=D_{T'}\) or \(W_{T'}\), equivariant nearby cycles are
\[
 R\Psi_{g,H}:\Dbc(Z_{\bar\eta'},E;H)
       \longrightarrow\Dbc(Z_{\bar t'},E;H),
 \qquad M\longmapsto L\bar i^*R\bar j_*M.
\]
If \(K_D\in\Dbc(Z,E)\), its geometric generic and special pullbacks carry
their canonical \(D_{T'}\)-descent data.  The cone of the specialization
map defines
\[
 R\Phi_{g,D}:\Dbc(Z,E)\longrightarrow
       \Dbc(Z_{\bar t'},E;D_{T'}).
\]
If \(K_W\in\DbcW(Z,E)\), the same cone formed with the strict Weil datum
defines
\[
 R\Phi_{g,W}:\DbcW(Z,E)\longrightarrow
       \Dbc(Z_{\bar t'},E;W_{T'}).
\]
In either case the defining distinguished triangle is
\[
 L\bar i^*K^{\mathrm{sh}}\longrightarrow
 R\Psi_{g,H}(L\bar j^*K^{\mathrm{sh}})\longrightarrow
 R\Phi_{g,H}(K)\longrightarrow L\bar i^*K^{\mathrm{sh}}[1],
\]
where \(H=D_{T'}\) for \(K_D\) and \(H=W_{T'}\) for \(K_W\).  Put
\(R^r\Phi_{g,H}=\mathcal H^rR\Phi_{g,H}\).
\end{definition}

\begin{theorem}[Restriction compatibility for cycles]
For \(M_D\in\Dbc(Z_{\bar\eta'},E;D_{T'})\) and
\(K_D\in\Dbc(Z,E)\), there are functorial isomorphisms
\[
 \begin{aligned}
 \operatorname{Res}^{D_{T'}}_{W_{T'}}R\Psi_{g,D_{T'}}(M_D)
   &\xrightarrow{\sim}
 R\Psi_{g,W_{T'}}\!\left(
   \operatorname{Res}^{D_{T'}}_{W_{T'}}M_D\right),\\
 \operatorname{Res}^{D_{T'}}_{W_{T'}}R\Phi_{g,D}(K_D)
   &\xrightarrow{\sim}
 R\Phi_{g,W}\!\left(\operatorname{For}^{D}_{W,Z}K_D\right).
 \end{aligned}
\]
They commute with the specialization triangles, cohomology sheaves, finite
coefficient extension, and finite residue-field extension.  The group
\(D_T\), or \(W_T\), acts on the local system in the first coordinate;
the distinct group \(D_{T'}\), or \(W_{T'}\), is the nearby-cycle action
coming from the second trait.
\end{theorem}

\begin{definition}[The three local transforms]
For \(V\in\Rep_E(W_T)\), let
\(V_{W,!}=j_{T!}\operatorname{Loc}_T^W(V)\), where
\(j_T:\eta\hookrightarrow T\).  Thus \(V_{W,!}\) is the extension by
zero of the strict lisse local Weil sheaf corresponding to \(V\).  The
degree-one Weil vanishing-cycle stalk at the geometric closed point
\(\bar t\) carries the action of \(W_{T'}\).  Define
\[
\begin{aligned}
 \mathcal F_\psi^{(0,\infty')}(V)
  &=\left(R^1\Phi_{g,W}
     (\operatorname{pr}^*V_{W,!}\otimes
      \overline{\mathcal L}_\psi(\pi/\pi'))\right)_{\bar t},\\
 \mathcal F_\psi^{(\infty,0')}(V)
  &=\left(R^1\Phi_{g,W}
     (\operatorname{pr}^*V_{W,!}\otimes
      \overline{\mathcal L}_\psi(\pi'/\pi))\right)_{\bar t},\\
 \mathcal F_\psi^{(\infty,\infty')}(V)
  &=\left(R^1\Phi_{g,W}
     (\operatorname{pr}^*V_{W,!}\otimes
      \overline{\mathcal L}_\psi(1/(\pi\pi')))\right)_{\bar t}.
\end{aligned}
\]
Thus each local transform is a functor
\[
       \mathcal F_\psi^{(\bullet,\bullet')}:
       \Rep_E(W_T)\longrightarrow\Rep_E(W_{T'}).
\]
The labels \(0\) and \(\infty\) record whether the local coordinate is a
coordinate or its inverse.  Interchanging \(T,\pi\) with \(T',\pi'\)
defines the three reverse transforms.
\end{definition}

\begin{theorem}[Compatibility with decomposition-group transforms]
For \(V\in\Rep_E(D_T)\), let
\(\mathcal F_{\psi,D}^{(\bullet,\bullet')}(V)\) denote the construction of
the preceding display using \(R\Phi_{g,D}\) and
\(j_{T!}\operatorname{Loc}_T^D(V)\).  It is an object of
\(\Rep_E(D_{T'})\), and there is a
functorial isomorphism in \(\Rep_E(W_{T'})\)
\[
 \operatorname{Res}^{D_{T'}}_{W_{T'}}
   \mathcal F_{\psi,D}^{(\bullet,\bullet')}(V)
 \xrightarrow{\sim}
 \mathcal F_\psi^{(\bullet,\bullet')}
   \bigl(\operatorname{Res}^{D_T}_{W_T}V\bigr).
\]
It is compatible with reverse transforms, changes of local coordinate,
finite coefficient extension, and finite residue-field extension.
\end{theorem}

\begin{theorem}[Global-to-local kernel comparison]
Embed the henselian traits at \(0\) and \(\infty\) in the two projective
lines used by the global Fourier transform.  Under the coordinates
\(x=\pi\), \(x^{-1}=\pi\), \(x'=\pi'\), and
\((x')^{-1}=\pi'\), restriction of the global kernel
\(\mathcal L_\psi(xx')\) identifies respectively with the three local
kernels having phases
\[
 \pi/\pi',\qquad \pi'/\pi,\qquad 1/(\pi\pi').
\]
For every \(V\in\Rep_E(W_T)\), compactification, proper base change,
and the nearby-cycle base-change morphism induce isomorphisms between the
vanishing cycles of the restricted global Fourier integrand and the three
local vanishing-cycle complexes in the preceding definition.  These
isomorphisms are equivariant for inertia and geometric Frobenius and are
compatible with changes of uniformizer, finite residue-field extension,
and finite coefficient extension.
\end{theorem}

\begin{theorem}[Degree-one concentration]
For every \(V\in\Rep_E(W_T)\), each of the three Weil vanishing-cycle complexes
used above is supported at the closed point of the special fiber and has
zero cohomology sheaves outside degree one.  Its degree-one stalk is a
finite-dimensional object of \(\Rep_E(W_{T'})\).  This statement
is functorial in \(V\) and is preserved by finite coefficient extension and
finite extension of the residue field.
\end{theorem}

\begin{theorem}[Exactness and base change of local Fourier transform]
The three degree-one constructions define exact functors
\[
 \mathcal F_\psi^{(0,\infty')},\quad
 \mathcal F_\psi^{(\infty,0')},\quad
 \mathcal F_\psi^{(\infty,\infty')}:
 \Rep_E(W_T)\longrightarrow\Rep_E(W_{T'}).
\]
They commute with finite direct sums and short exact sequences.  If
\(k_n/k\) is finite and
\(\psi_n=\psi\circ\operatorname{Tr}_{k_n/k}\), pullback of the traits to
\(k_n\) identifies their base changes with the three transforms defined by
\(\mathcal L_{\psi_n}\).  For every finite extension \(E'/E\),
\[
 \mathcal F_\psi^{(\bullet,\bullet')}(V)\otimes_EE'
 \xrightarrow{\sim}
 \mathcal F_\psi^{(\bullet,\bullet')}(V\otimes_EE').
\]
\end{theorem}

\begin{theorem}[The \((0,\infty')\) rank-and-break theorem]
The functor \(\mathcal F_\psi^{(0,\infty')}\) induces an exact equivalence
\[
 \Rep_E(W_T)\xrightarrow{\sim}\Rep_E(W_{T'})_{[0,1)}.
\]
For every \(V\),
\[
 \dim\mathcal F_\psi^{(0,\infty')}(V)=\dim V+\Sw(V),
 \qquad
 \Sw\bigl(\mathcal F_\psi^{(0,\infty')}(V)\bigr)=\Sw(V).
\]
If \(V\) is isoclinic of break \(\lambda\), its transform is isoclinic of
break \(\lambda/(1+\lambda)\), with multiplicity determined by the
displayed rank formula.
\end{theorem}

\begin{theorem}[The \((\infty,0')\) rank-and-break theorem]
The functor \(\mathcal F_\psi^{(\infty,0')}\) is zero on
\(\Rep_E(W_T)_{[1,\infty)}\) and induces an exact equivalence
\[
 \Rep_E(W_T)_{[0,1)}\xrightarrow{\sim}\Rep_E(W_{T'}).
\]
For \(V\in\Rep_E(W_T)_{[0,1)}\),
\[
 \dim\mathcal F_\psi^{(\infty,0')}(V)=\dim V-\Sw(V),
 \qquad
 \Sw\bigl(\mathcal F_\psi^{(\infty,0')}(V)\bigr)=\Sw(V).
\]
An isoclinic break \(0\le\lambda<1\) is carried to
\(\lambda/(1-\lambda)\).
\end{theorem}

\begin{theorem}[The \((\infty,\infty')\) rank-and-break theorem]
The functor \(\mathcal F_\psi^{(\infty,\infty')}\) is zero on
\(\Rep_E(W_T)_{[0,1]}\) and induces an exact equivalence
\[
 \Rep_E(W_T)_{(1,\infty)}
 \xrightarrow{\sim}\Rep_E(W_{T'})_{(1,\infty)}.
\]
For \(V\in\Rep_E(W_T)_{(1,\infty)}\),
\[
 \dim\mathcal F_\psi^{(\infty,\infty')}(V)=\Sw(V)-\dim V,
 \qquad
 \Sw\bigl(\mathcal F_\psi^{(\infty,\infty')}(V)\bigr)=\Sw(V).
\]
An isoclinic break \(\lambda>1\) is carried to
\(\lambda/(\lambda-1)\).
\end{theorem}

\begin{theorem}[Local Fourier inversion]
For the three equivalences, the reverse labels are respectively
\((\infty',0)\), \((0',\infty)\), and
\((\infty',\infty)\).  On the indicated slope categories there are natural
isomorphisms
\[
 \mathcal F_{\psi^{-1}}^{\mathrm{rev}}
 \mathcal F_\psi^{(\bullet,\bullet')}(V)\simeq V(-1),
 \qquad
 \mathcal F_\psi^{(\bullet,\bullet')}
 \mathcal F_{\psi^{-1}}^{\mathrm{rev}}(W)\simeq W(-1).
\]
Thus the reverse inverse-character transform followed by \((1)\) is a
quasi-inverse.  With the same character, the corresponding composites are
\[
 \mathcal F_\psi^{\mathrm{rev}}
 \mathcal F_\psi^{(\bullet,\bullet')}(V)
 \simeq a_T^*V(-1),
 \qquad
 \mathcal F_\psi^{(\bullet,\bullet')}
 \mathcal F_\psi^{\mathrm{rev}}(W)
 \simeq a_{T'}^*W(-1).
\]
The unit and counit are natural, coordinate-compatible, and satisfy the
two triangle identities.
\end{theorem}

\section{The type-3 transform at the origin}

Let \(\mathcal F_0\) be an arithmetically irreducible lisse Weil sheaf on
\(U_0\).  Assume that its geometric local system has no constant simple
subquotient and is unramified at infinity.  Put
\[
 \mathcal K_{\mathcal F,0}:=j_{0!*}^W(\mathcal F_0[1])
       \in\PervW(A_0,E),\qquad
 K_{\mathcal F}:=\operatorname{For}_A(\mathcal K_{\mathcal F,0})
       \simeq j_{!*}(\mathcal F[1])\simeq(j_*\mathcal F)[1].
\]
Let \(0'\in A'_0\) be the origin.  Put
\[
 U'_0=A'_0\setminus\{0'\},\qquad
 U'=U'_0\times_k\Fbar=A'\setminus\{0'\},
\]
let \(j'_0:U'_0\hookrightarrow A'_0\) be the open immersion, and let
\(j':U'\hookrightarrow A'\) be its geometric base change.  For chosen
geometric generic points, the decomposition-group construction for
\[
 T_{0'}^h=\operatorname{Spec}\mathcal O^h_{A'_0,0'}
 \quad\text{and}\quad
 T_\infty^h=\operatorname{Spec}\mathcal O^h_{\mathbf P^1_k,\infty}
\]
defines \(W_{0'},I_{0'}\) and \(W_\infty,I_\infty\).

\begin{theorem}[Type-3 Fourier sheaf]
There is an arithmetically irreducible lisse Weil sheaf
\(\mathcal F'_0\) on \(U'_0\) such that
\[
 \FT_{\psi,W}(\mathcal K_{\mathcal F,0})
       \xrightarrow{\sim}(j'_0)_{!*}^{W}(\mathcal F'_0[1])
\]
in \(\PervW(A'_0,E)\).  After applying \(\operatorname{For}_{A'}\), this
becomes
\[
 \FT_\psi(K_{\mathcal F})\simeq j'_{!*}(\mathcal F'[1])
                         \simeq(j'_*\mathcal F')[1].
\]
These are isomorphisms of whole perverse objects, so no punctual direct
summand is present.  Geometric irreducibility of the output is asserted only
when it is included in the geometric type-3 hypotheses.
\end{theorem}

Let
\[
 V'=\mathcal F'_{\bar\eta_{0'}}\in\Rep_E(W_{0'}),\qquad
 V_\infty=\mathcal F_{\bar\eta_\infty}\in\Rep_E(W_\infty).
\]
The rational points \(0'\) and \(\infty\) identify the residue Galois
quotients \(W_{0'}/I_{0'}\) and \(W_\infty/I_\infty\) with the infinite
cyclic subgroup generated by geometric Frobenius.  Thus an unramified Weil
representation at \(\infty\) can be regarded as one at \(0'\).

\begin{definition}[Stationary-phase morphisms]
Compact-support specialization at \(0'\), followed by the
global-to-local comparison for vanishing cycles, defines a canonical map
\[
 a_{\mathrm{sp}}:H_c^1(A,j_*\mathcal F)\longrightarrow V'.
\]
The Weil-group local Fourier functor
\(\mathcal F_\psi^{(\infty,0')}:
\Rep_E(W_\infty)_{[0,1)}\to\Rep_E(W_{0'})\), together with the
global-to-local comparison, defines a canonical map
\[
 b_{\mathrm{sp}}:V'\longrightarrow
 (V_\infty)^{I_\infty}(-1).
\]
The source of \(a_{\mathrm{sp}}\) and the target of
\(b_{\mathrm{sp}}\) are regarded as \(W_{0'}\)-representations through
the common quotient
\(W_{0'}/I_{0'}\simeq W_\infty/I_\infty\).
\end{definition}

\begin{proposition}[Stationary-phase equivariance and zero composite]
Both \(a_{\mathrm{sp}}\) and \(b_{\mathrm{sp}}\) are
\(W_{0'}\)-equivariant, including equivariance for geometric Frobenius,
and
\[
 b_{\mathrm{sp}}\circ a_{\mathrm{sp}}=0.
\]
The Frobenius on the last term is the Frobenius on
\((V_\infty)^{I_\infty}\) multiplied by \(q\), as prescribed by the
Tate twist \((-1)\).
\end{proposition}

\begin{proposition}[Left exactness of stationary phase]
The map
\(a_{\mathrm{sp}}:H_c^1(A,j_*\mathcal F)\to V'\) is injective.
\end{proposition}

\begin{proposition}[Middle exactness of stationary phase]
Inside \(V'\),
\[
 \operatorname{im}(a_{\mathrm{sp}})=\ker(b_{\mathrm{sp}}).
\]
\end{proposition}

\begin{proposition}[Right exactness of stationary phase]
The map
\(b_{\mathrm{sp}}:V'\to(V_\infty)^{I_\infty}(-1)\) is surjective.
\end{proposition}

\begin{theorem}[Local Fourier stationary-phase sequence]
The preceding maps form a canonical short exact sequence of continuous
\(W_{0'}\)-representations
\[
 0\longrightarrow H_c^1(A,j_*\mathcal F)
 \xrightarrow{a_{\mathrm{sp}}}V'
 \xrightarrow{b_{\mathrm{sp}}}(V_\infty)^{I_\infty}(-1)
 \longrightarrow0.
\]
It is natural in morphisms of type-3 input sheaves, independent of the
chosen local parameters up to the canonical coordinate-change
isomorphisms, and compatible with finite ground-field and coefficient-field
extensions.
\end{theorem}

\begin{theorem}[Stationary-phase inertia invariants]
The inclusion \(a_{\mathrm{sp}}\) factors through inertia invariants and
induces a Frobenius-equivariant isomorphism
\[
 H_c^1(A,j_*\mathcal F)\xrightarrow{\sim}(V')^{I_{0'}}.
\]
Consequently
\[
 W_1(V')=V'/(V')^{I_{0'}}
 \xrightarrow{\sim}(V_\infty)^{I_\infty}(-1).
\]
Inertia \(I_{0'}\) acts unipotently of echelon two on \(V'\).  Under the
standing unramifiedness assumption at infinity,
\((V_\infty)^{I_\infty}=V_\infty\).
\end{theorem}

\begin{proposition}[Normalization at the origin]
For a geometric point \(i_{0'}:\operatorname{Spec}\Fbar\to A'\), the
Fourier stalk formula and the chosen shifts give
\[
 H^{-1}\!\left(i_{0'}^*\FT_\psi(K_{\mathcal F})\right)
       \simeq H_c^1(A,j_*\mathcal F).
\]
Moreover,
\[
 H_c^0(A,j_*\mathcal F)=0,\qquad H_c^2(A,j_*\mathcal F)=0.
\]
\end{proposition}

\section{Realness and local weights}

\begin{definition}[Constant Weil character]
For \(c\in E^\times\), let \(\chi_c\) be the rank-one constant lisse Weil
sheaf on a geometrically connected \(k\)-scheme on which geometric
Frobenius over \(\mathbf F_{q^n}\) acts by \(c^n\).  It is
\(\iota\)-pure of the unique real weight \(t\) satisfying
\[
                         |\iota(c)|=q^{t/2}.
\]
\end{definition}

\begin{theorem}[Rank-one class-field decomposition]
Let \(C_0\) be a smooth geometrically connected curve over \(k\), and let
\(\mathcal L_0\) be a rank-one lisse Weil sheaf.  After a finite extension
of the coefficient field, there are a finite-order rank-one lisse Weil
sheaf \(\mathcal L_{\mathrm{fin},0}\) and a scalar \(c\in E^\times\) such
that
\[
 \mathcal L_0\simeq\mathcal L_{\mathrm{fin},0}\otimes\chi_c.
\]
Equivalently, for every closed point \(x\), its Frobenius eigenvalue has the
form
\[
 \alpha_x=\varepsilon_xc^{\deg(x)},
 \qquad \varepsilon_x\text{ a root of unity}.
\]
The finite character is obtained from degree-zero divisor classes through
global class-field reciprocity; it is unchanged after replacing the curve
by a dense open on which \(\mathcal L_0\) is lisse.
\end{theorem}

\begin{definition}[Determinantal weight]
Let \(\mathcal S_0\) be arithmetically irreducible of rank \(r>0\).  Apply
the rank-one decomposition to
\(\det\mathcal S_0\simeq\mathcal L_{\mathrm{fin},0}\otimes\chi_c\).
Its \(\iota\)-determinantal weight is
\[
 d=\frac{2\log|\iota(c)|}{r\log q}.
\]
Thus, if \(\alpha_{x,1},\ldots,\alpha_{x,r}\) are the Frobenius
eigenvalues at a degree-\(e\) point, then
\[
 \left|\iota\!\left(\prod_{j=1}^r\alpha_{x,j}\right)\right|
       =q^{erd/2}.
\]
\end{definition}

\begin{definition}[Global Weil group and degree cosets]
Put \(C=C_0\times_k\Fbar\), choose \(\bar u\to C\), and let
\(\langle\Frob_q\rangle\simeq\mathbf Z\) be the subgroup of
\(\operatorname{Gal}(\Fbar/k)\) generated by geometric Frobenius.  The
global Weil group is the fiber product
\[
 W(C_0,\bar u)=
 \pi_1(C_0,\bar u)
 \mathop{\times}_{\operatorname{Gal}(\Fbar/k)}
 \langle\Frob_q\rangle .
\]
Its degree homomorphism and degree cosets are
\[
 \deg:W(C_0,\bar u)\longrightarrow\mathbf Z,
 \qquad
 W(C_0,\bar u)^{(m)}=\deg^{-1}(m).
\]
The kernel \(W(C_0,\bar u)^{(0)}=\pi_1(C,\bar u)\) is open with its
profinite topology; every degree coset has the translated topology.  A
strict lisse Weil sheaf \(\mathcal S_0\) on \(C_0\) is equivalently a
finite-dimensional continuous representation
\[
 \rho_W:W(C_0,\bar u)\longrightarrow\operatorname{GL}(V),
 \qquad V=(\mathcal S_0)_{\bar u}.
\]
Under this equivalence the operator attached to a closed point of degree
\(e\) is the action of a conjugacy class contained in
\(W(C_0,\bar u)^{(e)}\).
\end{definition}

\begin{definition}[Arithmetic and geometric monodromy groups]
For the global Weil representation \(\rho_W\) of \(\mathcal S_0\), put
\[
 \begin{aligned}
 G&=\overline{\rho_W(W(C_0,\bar u))}^{\mathrm{Zar}},\\
 G_{\mathrm{geom}}
   &=\overline{\rho_W(W(C_0,\bar u)^{(0)})}^{\mathrm{Zar}},\\
 G^{(m)}&=
   \overline{\rho_W(W(C_0,\bar u)^{(m)})}^{\mathrm{Zar}}.
 \end{aligned}
\]
For every \(m\), the locally closed subset \(G^{(m)}\) is a left and
right coset of \(G_{\mathrm{geom}}\) in \(G\).  An element of
\(G^{(m)}(\overline E)\) is, by definition, an algebraic-monodromy element
of degree \(m\).
Let \(R_u(G_{\mathrm{geom}}^\circ)\) be the unipotent radical and set
\[
 G_{\mathrm{red}}=G_{\mathrm{geom}}^\circ/
                   R_u(G_{\mathrm{geom}}^\circ).
\]
\end{definition}

\begin{theorem}[Reductive monodromy comparison]
The monodromy image is Zariski dense in \(G\).  The unipotent radical of
\(G_{\mathrm{geom}}^\circ\) acts trivially on the semisimplification of
every tensor or exterior-power representation generated by \(V\).  Hence
these semisimplifications factor through \(G_{\mathrm{red}}\), and the
factorization maps admit equivariant inverses onto their images.
\end{theorem}

\begin{theorem}[Central positive-degree weight detector]
After a finite extension of \(k\), there are an integer \(m>0\) and an
algebraic-monodromy element \(z\in G^{(m)}(\overline E)\) whose action
centralizes \(G_{\mathrm{red}}\).  If \(W\) is an
irreducible \(G_{\mathrm{red}}\)-constituent of determinantal weight
\(d_W\), every eigenvalue \(\zeta\) of \(z\) on \(W\) satisfies
\[
                         |\iota(\zeta)|=q^{md_W/2}.
\]
This applies to every constituent selected in tensor and exterior powers of
\(V\).
\end{theorem}

\begin{definition}[Numerical weight of a local eigenvalue]
For a degree-\(e\) point and a nonzero Frobenius eigenvalue \(\alpha\), put
\[
 \nu_\iota(e,\alpha)=
       \frac{2\log|\iota(\alpha)|}{\log(q^e)}.
\]
Thus \(\alpha\) is pure of weight \(\nu_\iota(e,\alpha)\).
\end{definition}

\begin{theorem}[Exterior-power constituent selection]
For \(1\le s\le r\) and a product
\(\alpha_{x,i_1}\cdots\alpha_{x,i_s}\) of distinct local eigenvalue
occurrences, there is an irreducible monodromy constituent
\(\mathcal T_0\) of \((\bigwedge^s\mathcal S_0)^{\mathrm{ss}}\), together with
equivariant maps
\[
 i_{\mathcal T}:\mathcal T_0\longrightarrow
       (\bigwedge^s\mathcal S_0)^{\mathrm{ss}},
 \qquad
 p_{\mathcal T}:(\bigwedge^s\mathcal S_0)^{\mathrm{ss}}
       \longrightarrow\mathcal T_0,
\]
such that
\[
 p_{\mathcal T}i_{\mathcal T}=1_{\mathcal T_0},\qquad
 (i_{\mathcal T}p_{\mathcal T})^2=i_{\mathcal T}p_{\mathcal T},
\]
and the selected product occurs in the local spectrum of \(\mathcal T_0\).
The constituent and maps are compatible with the central weight detector
and finite coefficient extension.
\end{theorem}

\begin{theorem}[Low-degree cohomology on a curve]
Let \(j:C_0\hookrightarrow\overline C_0\) be a smooth compactification and
let \(\mathcal T_0\) be any irreducible tensor- or exterior-power
constituent above.  Assume that \(C_0\) is nonproper.  There are
Frobenius-equivariant identifications
\[
 H^0(C,\mathcal T)=V_{\mathcal T}^{\pi_1(C)},\qquad
 H_c^0(C,\mathcal T)=0,\qquad
 H_c^2(C,\mathcal T)\simeq
 H^0(C,\mathcal T^\vee)^\vee(-1).
\]
The central weight detector determines the Frobenius moduli on the two
nonzero outer terms, and all three identifications commute with finite
ground- and coefficient-field extension.
\end{theorem}

\begin{definition}[Euler and cohomological \(L\)-functions]
For a lisse Weil sheaf \(\mathcal T_0\) on \(C_0\), define
\[
 L_{\mathrm{Euler}}(C_0,\mathcal T_0,T)=
 \prod_{x\in|C_0|}
 \det(1-T^{\deg(x)}\Frob_x\mid\mathcal T_{\bar x})^{-1}
\]
as a formal power series, and define
\[
 L_{\mathrm{coh}}(C_0,\mathcal T_0,T)=
 \prod_{i=0}^2
 \det(1-T\Frob_q\mid H_c^i(C,\mathcal T))^{(-1)^{i+1}}
\]
as a rational function.
\end{definition}

\begin{theorem}[Trace--\(L\)-function comparison]
The Grothendieck--Lefschetz trace formula over every finite extension of
\(k\) gives an equality in \(E[[T]]\)
\[
 L_{\mathrm{Euler}}(C_0,\mathcal T_0,T)
       =L_{\mathrm{coh}}(C_0,\mathcal T_0,T).
\]
The coefficient of \(T^n/n\) in the logarithm of either side is
\[
 \sum_{x\in C_0(\mathbf F_{q^n})}
       \operatorname{Tr}(F_{x,n}\mid\mathcal T_{\bar x}).
\]
\end{theorem}

\begin{theorem}[Euler-product convergence]
Let
\[
 M_{\mathcal T}=\sup_{x,\alpha}
       |\iota(\alpha)|^{1/\deg(x)},
\]
where \(\alpha\) runs through the local Frobenius eigenvalues of
\(\mathcal T_0\).  This supremum is finite.  The point-count bounds obtained from a finite
quasi-finite stratification of \(C_0\) imply absolute convergence of the
logarithmic Euler product on
\[
                         |T|<(qM_{\mathcal T})^{-1},
\]
uniformly on every smaller closed disk.
\end{theorem}

\begin{theorem}[Paired tensor-factor reality]
Let \(\mathcal S_0\) be an irreducible constituent of an \(\iota\)-real
lisse Weil sheaf \(\mathcal G_0\), and let \(\mathcal T_0\) be selected
from a tensor or exterior power of \(\mathcal S_0\).  After a finite
extension of \(E\), there is a conjugate constituent
\(\mathcal T_0^c\), with Frobenius-equivariant inclusion and retraction
into the corresponding semisimple tensor construction \(\mathcal A_0\),
\[
 i_c:\mathcal T_0^c\longrightarrow\mathcal A_0,
 \qquad p_c:\mathcal A_0\longrightarrow\mathcal T_0^c,
 \qquad p_ci_c=1,\qquad(i_cp_c)^2=i_cp_c.
\]
For every closed point \(x\), the \(\iota\)-image of the Frobenius
spectrum on \((\mathcal T_0^c)_{\bar x}\) is the complex conjugate, with
multiplicity, of the \(\iota\)-image of the spectrum on
\((\mathcal T_0)_{\bar x}\).

Put
\[
             \mathcal R_0=\mathcal T_0\otimes_E\mathcal T_0^c .
\]
This is the \emph{paired correlation factor}.  Its local polynomial is
\[
 P_{\mathcal R,x}(U)=
 \det(1-U\Frob_x\mid\mathcal R_{\bar x}).
\]
For every closed \(x\),
\(\iota(P_{\mathcal R,x})\in\mathbf R[U]\); and for every \(n\ge1\) and
every \(x\in C_0(\mathbf F_{q^n})\),
\[
 \iota\!\left(\operatorname{Tr}
       (F_{x,n}\mid\mathcal R_{\bar x})\right)
 =\left|\iota\!\left(\operatorname{Tr}
       (F_{x,n}\mid\mathcal T_{\bar x})\right)\right|^2.
\]
More explicitly, put
\[
 i_{\mathcal R}=1_{\mathcal T_0}\otimes i_c,\qquad
 p_{\mathcal R}=1_{\mathcal T_0}\otimes p_c,\qquad
 e_{\mathcal R}=i_{\mathcal R}p_{\mathcal R}.
\]
Then \(i_{\mathcal R}\) and \(p_{\mathcal R}\) exhibit
\(\mathcal R_0\) as a Frobenius-equivariant direct summand of
\(\mathcal T_0\otimes_E\mathcal A_0\), and
\[
 p_{\mathcal R}i_{\mathcal R}=1_{\mathcal R_0},\qquad
 e_{\mathcal R}^2=e_{\mathcal R}.
\]
All these data commute with finite coefficient extension.
\end{theorem}

\begin{definition}[Cohomological boundary-pole datum]
For the paired correlation factor \(\mathcal R_0\), set
\[
 \begin{aligned}
 N_{\mathcal R}(T)
   &=\iota\!\left(\det(1-T\Frob_q\mid H_c^1(C,\mathcal R))\right),\\
 D_{\mathcal R}(T)
   &=\iota\!\left(\det(1-T\Frob_q\mid H_c^0(C,\mathcal R))\right)
     \,\iota\!\left(\det(1-T\Frob_q\mid H_c^2(C,\mathcal R))\right).
 \end{aligned}
\]
Thus
\[
 \iota L_{\mathrm{coh}}(C_0,\mathcal R_0,T)
       =\frac{N_{\mathcal R}(T)}{D_{\mathcal R}(T)}.
\]
Let \(H_{\mathcal R}\in\mathbf R[T]\) be the common divisor normalized by
\(H_{\mathcal R}(0)=1\) such that
\[
 P_{\mathcal R}=N_{\mathcal R}/H_{\mathcal R},\qquad
 Q_{\mathcal R}=D_{\mathcal R}/H_{\mathcal R}
\]
are coprime and have constant term one.  Write
\[
 \frac{P_{\mathcal R}(T)}{Q_{\mathcal R}(T)}
      =\sum_{n\ge0}a_nT^n
\]
at the origin, let \(r_{\mathcal R}\) be the radius of convergence of this
Taylor series, and call
\[
             \mathcal C_{\mathcal R}
             =\{T\in\mathbf C:|T|=r_{\mathcal R}\}
\]
its boundary circle.  For \(\lambda\in\mathbf C^\times\), put
\[
 m_D(\lambda)=\operatorname{ord}_{T=\lambda^{-1}}D_{\mathcal R}(T),
 \qquad
 m_N(\lambda)=\operatorname{ord}_{T=\lambda^{-1}}N_{\mathcal R}(T).
\]
It is a \emph{boundary-extremal uncancelled denominator eigenvalue} when
\[
 |\lambda|=r_{\mathcal R}^{-1},\qquad
 m_D(\lambda)>m_N(\lambda),
\]
or, equivalently,
\(Q_{\mathcal R}(\lambda^{-1})=0\) and
\(P_{\mathcal R}(\lambda^{-1})\ne0\).
\end{definition}

\begin{theorem}[Trace--Euler pole detection]
Let \(\mathcal S_0\) be an irreducible constituent of an \(\iota\)-real
lisse Weil sheaf \(\mathcal G_0\), and let \(\mathcal T_0\) be a constituent
selected from a tensor or exterior power of \(\mathcal S_0\), with
determinantal weight \(d_{\mathcal T}\).  Form \(\mathcal R_0\),
\(N_{\mathcal R}\), \(D_{\mathcal R}\), and \(r_{\mathcal R}\) as above.
Then \(0<r_{\mathcal R}<\infty\), the Taylor coefficients \(a_n\) are
nonnegative real numbers, and
\[
 \lambda_{\mathcal R}=r_{\mathcal R}^{-1}\in\mathbf R_{>0}
\]
is a boundary-extremal uncancelled denominator eigenvalue.  Explicitly,
\[
 Q_{\mathcal R}(r_{\mathcal R})=0,\qquad
 P_{\mathcal R}(r_{\mathcal R})\ne0,\qquad
 m_D(\lambda_{\mathcal R})>m_N(\lambda_{\mathcal R}).
\]
For every closed point \(x\in|C_0|\) and every eigenvalue
\(\beta\in\overline E^\times\) of \(\Frob_x\) on
\((\mathcal T_0)_{\bar x}\), this pole and the central degree detector give
\[
 |\iota(\beta)|\le q_x^{(d_{\mathcal T}+1)/2},
 \qquad
 \nu_\iota(\deg(x),\beta)\le d_{\mathcal T}+1.
\]
For every tensor power, finite extension of the ground field, or finite
extension of the coefficient field, the corresponding paired factor and
cohomological boundary-pole datum satisfy the same displayed multiplicity
inequality and local-eigenvalue bound.
\end{theorem}

\begin{theorem}[Tensor-power limiting inequality]
Let \(\mathcal S_0\) be an irreducible constituent of an \(\iota\)-real
lisse Weil sheaf, let \(d\) be its determinantal weight, and let \(\alpha\)
be one of its local eigenvalues at a degree-\(e\) point.  For every integer
\(k\ge1\), an irreducible constituent of
\((\mathcal S_0^{\otimes k})^{\mathrm{ss}}\) contains the occurrence
\(\alpha^k\), has
central determinantal weight \(kd\), and satisfies
\[
                         \nu_\iota(e,\alpha)\le d+\frac1k.
\]
\end{theorem}

\begin{corollary}[Removal of the tensor-power error]
Under the hypotheses of the preceding theorem, every local eigenvalue
\(\alpha\) of \(\mathcal S_0\) satisfies
\[
                         \nu_\iota(e,\alpha)\le d.
\]
\end{corollary}

\begin{theorem}[Determinant numerical-weight identity]
Let \(\mathcal S_0\) be an arithmetically irreducible lisse Weil sheaf of
rank \(r>0\), and let \(d\) be its \(\iota\)-determinantal weight.  If
\(x\) is a closed point of degree \(e\) and
\(\alpha_{x,1},\ldots,\alpha_{x,r}\) is the Frobenius spectrum on
\((\mathcal S_0)_{\bar x}\), then
\[
             \sum_{j=1}^r\nu_\iota(e,\alpha_{x,j})=rd.
\]
\deps{The determinant-product formula in the definition of determinantal
weight and the definition of \(\nu_\iota(e,-)\).}
\end{theorem}

\begin{lemma}[Finite upper-bound saturation]
If \(a_1,\ldots,a_r\le d\) and
\(a_1+\cdots+a_r=rd\), then \(a_j=d\) for every \(j\).
\end{lemma}

\begin{theorem}[Deligne's PUR endpoint]
Let \(\mathcal S_0\) be an arithmetically irreducible constituent of an
\(\iota\)-real lisse Weil sheaf, let \(\mathcal S\) be its geometric
pullback, and let \(d\) be its determinantal weight.  Then
\[
 \nu_\iota(\deg(x),\alpha)=d
\]
for every closed point \(x\) and every Frobenius eigenvalue \(\alpha\) on
\(\mathcal S_{\bar x}\).  Thus \(\mathcal S_0\) is punctually
\(\iota\)-pure of weight \(d\).
\deps{Rank-one class-field decomposition, reductive monodromy, the central
positive-degree weight detector, low-degree cohomology, the trace--\(L\)
identity, Euler-product convergence, pole detection, tensor-power limiting,
the determinant numerical-weight identity, finite upper-bound saturation,
and exterior-power constituent selection.}
\end{theorem}

\begin{corollary}[Finite-determinant purity]
If, in addition, \(\det(\mathcal S_0)\) has finite order, then its
determinantal weight is zero and \(\mathcal S_0\) is punctually
\(\iota\)-pure of weight zero.
\end{corollary}

\begin{proposition}[Constant-character weight translation]
If \(\mathcal S_0\otimes\chi_c^{-1}\) is punctually pure of weight zero and
\(|\iota(c)|=q^{t/2}\), then \(\mathcal S_0\) is punctually pure of weight
\(t\).  The assertion is unchanged by finite coefficient extension.
\end{proposition}

\begin{theorem}[Deligne's realness criterion]
Let \(C_0\) be a smooth geometrically connected curve over \(k\).  Every
arithmetically irreducible lisse subquotient of an \(\iota\)-real lisse Weil
sheaf on \(C_0\) has a determinantal weight and is punctually
\(\iota\)-pure of that real weight.
\end{theorem}

\begin{theorem}[Ordinary-extension boundary stalk]
Let \(j:U\hookrightarrow C\) be a dense open immersion of smooth curves and
let \(\mathcal G\) be a lisse Weil sheaf on \(U\), punctually
\(\iota\)-pure of weight \(v\in\mathbf R\).  For a missing closed point
\(s\) and \(V=\mathcal G_{\bar\eta_s}\), there is a canonical
Frobenius-equivariant isomorphism
\[
 (j_*\mathcal G)_{\bar s}\xrightarrow{\sim}V^{I_s}.
\]
\end{theorem}

\begin{theorem}[Boundary-invariant upper weights]
With the preceding notation, the Frobenius module \(V^{I_s}\) has weights
at most \(v\).
\end{theorem}

\begin{theorem}[First boundary direct-image stalk]
There is a canonical Frobenius-equivariant isomorphism
\[
 (R^1j_*\mathcal G)_{\bar s}\xrightarrow{\sim}H^1(I_s,V),
\]
where continuous group cohomology is taken after restriction of \(V\) to
\(I_s\), and geometric Frobenius acts through conjugation on \(I_s\) and
through its given action on \(V\).
\end{theorem}

\begin{theorem}[Local-duality boundary pairing]
Local duality gives a Frobenius-equivariant perfect pairing, equivalently a
canonical isomorphism,
\[
 H^1(I_s,V)^\vee\xrightarrow{\sim}(V^\vee(1))^{I_s}.
\]
The right-hand side has weights at most \(-v-2\); hence
\(H^1(I_s,V)\) has weights at least \(v+2\).
\end{theorem}

\begin{corollary}[Boundary weights]
At every missing closed point \(s\),
\((j_*\mathcal G)_{\bar s}\) has weights at most \(v\), while
\((R^1j_*\mathcal G)_{\bar s}\) has weights at least \(v+2\).  Explicitly,
every eigenvalue \(\beta\) on the latter space satisfies
\[
                 |\iota(\beta)|\ge q_s^{(v+2)/2}.
\]
\end{corollary}

\begin{theorem}[Unramified extension across the boundary]
If \(I_s\) acts trivially on \(V\), then \(\mathcal G\) extends uniquely as
a lisse Weil sheaf across \(s\), and its fiber at \(s\) is punctually pure
of weight \(v\).
\end{theorem}

\begin{theorem}[Echelon-two upper bound]
In the preceding setting, assume that \(I_s\) acts trivially on
\(W_1(\mathcal G_{\bar\eta_s})\).  Every Frobenius eigenvalue \(\alpha\)
on this quotient satisfies
\[
             |\iota(\alpha)|\le q_s^{(v+1)/2}.
\]
This is an upper bound and is not by itself a purity statement.
\end{theorem}

\begin{definition}[Local characteristic polynomial and Euler factor]
For a constructible Weil sheaf \(\mathcal H_0\) on a curve and a closed
point \(x\), put
\[
 P_x(\mathcal H_0,T)=
 \det(1-T\Frob_x\mid\mathcal H_{\bar x}),\qquad
 L_x(\mathcal H_0,T)=P_x(\mathcal H_0,T^{\deg(x)})^{-1}.
\]
The first polynomial records the local Frobenius module; the second is the
corresponding factor in the Euler product with the convention used here.
\end{definition}

\begin{lemma}[Newton reconstruction and the reality criterion]
Let \(V\) be a finite-dimensional \(E\)-vector space of dimension \(d\),
let \(F\in\operatorname{End}_E(V)\), and write
\[
 \det(1-TF\mid V)=\sum_{r=0}^{d}c_rT^r,
 \qquad c_0=1.
\]
The coefficients are uniquely specified by the Newton recursions
\[
 r c_r+\sum_{m=1}^{r}c_{r-m}\operatorname{Tr}(F^m\mid V)=0
 \qquad(1\leq r\leq d).
\]
The following conditions are equivalent:
\begin{enumerate}
\item \(\iota(\det(1-TF\mid V))\in\mathbf R[T]\);
\item \(\iota(\operatorname{Tr}(F^m\mid V))\in\mathbf R\) for
      \(1\leq m\leq d\);
\item \(\iota(\operatorname{Tr}(F^m\mid V))\in\mathbf R\) for every
      \(m\geq1\).
\end{enumerate}
\end{lemma}

\begin{definition}[Boundary-real conjugate companion]
Let \(\mathcal F_0^c\) be an \(\iota\)-conjugate companion of
\(\mathcal F_0\), assumed arithmetically irreducible, unramified at
infinity, and without a constant geometric simple subquotient.  Thus it
satisfies the same type-3 hypotheses as \(\mathcal F_0\).  The ordinary
direct image equipped with its induced geometric-Frobenius descent datum is
the constructible Weil sheaf
\[
       \mathcal M_0=j_{0*}(\mathcal F_0\oplus\mathcal F_0^c)
       \quad\text{on }A_0.
\]
The companion is \emph{boundary-real} if
\[
                 \iota(P_x(\mathcal M_0,T))\in\mathbf R[T]
\]
for every closed point \(x\in A_0\), including the points of
\(A_0\setminus U_0\).
\end{definition}

\begin{theorem}[Paired-transform concentration and lissity]
Assume the type-3 hypotheses and let \(\mathcal M_0\) be as above.  For
\(\chi\in\{\psi,\psi^{-1}\}\), the restricted strict Weil complex
\[
 C_{\chi,0}=(j'_0)^*\FT_{\chi,W}(\mathcal M_0[1])
       \in\DbcW(U'_0,E)
\]
satisfies
\[
 \mathcal H^i(C_{\chi,0})=0\quad(i\ne-1),
 \qquad
 \mathcal H^{-1}(C_{\chi,0})
 \text{ is a finite-rank lisse Weil sheaf on }U'_0.
\]
The concentration and lissity statements commute with the forgetful
functor, finite coefficient extension, and finite extension of the ground
field.
\end{theorem}

\begin{definition}[Paired transformed lisse sheaf]
Define the lisse Weil sheaf \(\mathcal G'_0\) on \(U'_0\) by
\[
 \mathcal G'_0=\mathcal H^{-1}(j'_0)^*\!\left(
  \FT_{\psi,W}(\mathcal M_0[1])\oplus
  \FT_{\psi^{-1},W}(\mathcal M_0[1])
 \right).
\]
Equivalently, the preceding theorem supplies the unique isomorphism
\[
 \mathcal G'_0[1]\xrightarrow{\sim}
 (j'_0)^*\!\left(
  \FT_{\psi,W}(\mathcal M_0[1])\oplus
  \FT_{\psi^{-1},W}(\mathcal M_0[1])
 \right).
\]
\end{definition}

\begin{theorem}[Paired Fourier power-trace identity]
Let \(\mathcal M_0=j_{0*}(\mathcal F_0\oplus\mathcal F_0^c)\) be
boundary-real.  For every \(n\geq1\), the degree-\(n\) trace functions
\[
 t_{\FT_{\psi,W}(\mathcal M_0[1])}^{(n)},\quad
 t_{\FT_{\psi^{-1},W}(\mathcal M_0[1])}^{(n)}:
 A'_0(\mathbf F_{q^n})\longrightarrow E
\]
of the two normalized Fourier transforms satisfy, for every
\(x'\in A'_0(\mathbf F_{q^n})\),
\[
 \begin{aligned}
 &t_{\FT_{\psi,W}(\mathcal M_0[1])}^{(n)}(x')
 +t_{\FT_{\psi^{-1},W}(\mathcal M_0[1])}^{(n)}(x')\\
 &\quad=\sum_{x\in A_0(\mathbf F_{q^n})}t_{\mathcal M_0}^{(n)}(x)
 \left[
  \psi\!\left(\operatorname{Tr}_{\mathbf F_{q^n}/k}(xx')\right)
 +\psi^{-1}\!\left(\operatorname{Tr}_{\mathbf F_{q^n}/k}(xx')\right)
 \right].
 \end{aligned}
\]
For every \(n\geq1\) and every \(x'\in A'_0(\mathbf F_{q^n})\), the image
under \(\iota\) of the right-hand side is real.  For every closed point
\(x'\in U'_0\) and every \(m\geq1\),
\[
 \iota\!\left(
  \operatorname{Tr}\!\left(\Frob_{x'}^m\mid
  (\mathcal G'_0)_{\bar x'}\right)
 \right)\in\mathbf R.
\]
\end{theorem}

\begin{corollary}[Reality of the paired transformed local factor]
For every closed point \(x'\in U'_0\),
\[
             \iota(P_{x'}(\mathcal G'_0,T))\in\mathbf R[T].
\]
This is a statement about individual local factors and asserts no equality
of global Euler products.
\end{corollary}

\begin{theorem}[Paired Fourier real envelope]
Assume that \(\mathcal F_0^c\) is a boundary-real conjugate companion of the
type-3 sheaf \(\mathcal F_0\), and let \(\mathcal G'_0\) be the paired
transformed lisse sheaf just defined.  Then \(\mathcal G'_0\) is an
\(\iota\)-real Weil sheaf.  There are morphisms of lisse strict Weil sheaves
on \(U'_0\)
\[
 i':\mathcal F'_0\longrightarrow\mathcal G'_0,
 \qquad
 r':\mathcal G'_0\longrightarrow\mathcal F'_0
\]
such that
\[
 r'i'=1_{\mathcal F'_0}.
\]
Thus the strict-Weil endomorphism
\[
 e':=i'r'\in
 \operatorname{End}_{\operatorname{LisseW}(U'_0,E)}(\mathcal G'_0)
\]
satisfies
\[
 (e')^2=e',\qquad
 \operatorname{im}(e')\xrightarrow{\,r'\,}\mathcal F'_0,
 \quad
 \mathcal F'_0\xrightarrow{\,i'\,}\operatorname{im}(e')
\]
with both displayed maps inverse to one another.  In particular,
\(\mathcal F'_0\) is an arithmetic direct summand of \(\mathcal G'_0\).
Equivalently,
\[
               \iota(P_{x'}(\mathcal G'_0,T))\in\mathbf R[T]
\]
for every closed point \(x'\in U'_0\).  The assertion is compatible with
every finite extension of \(k\); its local-factor equality is the equality
of the characteristic polynomials recovered from all Frobenius power traces,
not merely an equality of one global Euler product.
\end{theorem}

\begin{corollary}[Purity of the Fourier output]
Under the hypotheses of the paired Fourier real-envelope theorem, the
arithmetically irreducible sheaf \(\mathcal F'_0\) is punctually
\(\iota\)-pure of a uniquely determined real weight \(w'\).
\end{corollary}

\begin{theorem}[Integral-weight conjugate companion on the open curve]
If \(\mathcal F_0\) is punctually \(\iota\)-pure of an integer weight
\(w\), then
\[
                     \mathcal F_0^c=\mathcal F_0^\vee(-w)
\]
is an \(\iota\)-conjugate companion on \(U_0\).  Under the type-3
hypotheses, it is arithmetically irreducible, unramified at infinity, and
has no constant geometric simple subquotient, so it again satisfies the
type-3 hypotheses.
\end{theorem}

\begin{definition}[Local inertia coinvariants]
For a local Weil representation \(V\in\Rep_E(W_s)\), define
\[
 V_{I_s}=V/\langle\sigma v-v:\sigma\in I_s,\ v\in V\rangle_E.
\]
It is a finite-dimensional representation of \(W_s/I_s\), with the
induced geometric-Frobenius operator.
\end{definition}

\begin{theorem}[Dual invariants and coinvariants]
For every \(V\in\Rep_E(W_s)\), evaluation induces a canonical
Frobenius-equivariant isomorphism
\[
 (V^\vee)^{I_s}\xrightarrow{\sim}(V_{I_s})^\vee.
\]
After a Tate twist, the same isomorphism identifies
\((V^\vee(-w))^{I_s}\) with \((V_{I_s})^\vee(-w)\).
\end{theorem}

\begin{theorem}[Boundary invariant--coinvariant symmetry]
Let \(\mathcal F_0\) be arithmetically irreducible and punctually
\(\iota\)-pure of an integer weight \(w\), and put
\(V=\mathcal F_{\bar\eta_s}\).  The Frobenius eigenvalue multiset on
\((V_{I_s})^\vee(-w)\), after applying \(\iota\), is the complex-conjugate
multiset of the Frobenius eigenvalues on \(V^{I_s}\).  This comparison is
compatible with finite coefficient and residue-field extension.
\end{theorem}

\begin{corollary}[Boundary characteristic-polynomial identity]
With \(\mathcal F_0^c=\mathcal F_0^\vee(-w)\), the preceding two theorems
give, for every \(s\in A_0\setminus U_0\),
\[
 \iota\det\!\left(1-T\Frob_s\mid
       (\mathcal F^c_{\bar\eta_s})^{I_s}\right)
 =
 \overline{\iota\det\!\left(1-T\Frob_s\mid
       \mathcal F_{\bar\eta_s}^{I_s}\right)}.
\]
\end{corollary}

\begin{theorem}[Deligne's boundary-polynomial theorem]
Let \(\mathcal F_0\) be arithmetically irreducible, punctually
\(\iota\)-pure of an integer weight \(w\), and put
\(\mathcal F_0^c=\mathcal F_0^\vee(-w)\).  For every
\(s\in A_0\setminus U_0\), the ordinary-extension stalks satisfy
\[
 \iota\det\!\left(1-T\Frob_s\mid
       (\mathcal F^c_{\bar\eta_s})^{I_s}\right)
 =
 \overline{\iota\det\!\left(1-T\Frob_s\mid
       \mathcal F_{\bar\eta_s}^{I_s}\right)}.
\]
Consequently \(\mathcal F_0^c\) is boundary-real in the sense defined
above.  The local comparison is through
\((V^\vee)^{I_s}\simeq(V_{I_s})^\vee\), not through the generally false
identity \((V^\vee)^{I_s}=(V^{I_s})^\vee\).
\end{theorem}

\begin{corollary}[Integral-weight Fourier purity]
If \(\mathcal F_0\) satisfies the type-3 Fourier-sheaf hypotheses and is
punctually \(\iota\)-pure of an integer weight, then the paired Fourier
sheaf \(\mathcal G'_0\) is \(\iota\)-real and its arithmetic direct summand
\(\mathcal F'_0\) is punctually \(\iota\)-pure of some real weight.
\end{corollary}

\begin{theorem}[Type-3 tame unipotence]
For the local representation
\(V'=\mathcal F'_{\bar\eta_{0'}}\in\Rep_E(W_{0'})\), inertia acts
trivially on \(W_1(V')=V'/(V')^{I_{0'}}\) and
\[
 (\sigma-1)V'\subset(V')^{I_{0'}},\qquad
 (\sigma-1)^2=0\quad(\sigma\in I_{0'}).
\]
Wild inertia acts trivially, and the remaining inertia action factors
through the maximal pro-\(\ell\) quotient of tame inertia.
\end{theorem}

\begin{definition}[Tame character at \(0'\)]
Fix a compatible system of \(\ell\)-power roots of a local parameter.  The
associated tame character is
\[
 t_\ell:I_{0'}\longrightarrow\mathbf Z_\ell(1).
\]
For a chosen lift \(\Phi_{0'}\in W_{0'}\) of geometric Frobenius, it is
normalized by
\[
 t_\ell(\Phi_{0'}\sigma\Phi_{0'}^{-1})
       =q^{-1}t_\ell(\sigma).
\]
\end{definition}

\begin{theorem}[Type-3 monodromy logarithm]
There is a unique \(E\)-linear morphism
\[
 N:V'\longrightarrow V'(-1)
\]
such that \(\rho(\sigma)=1+t_\ell(\sigma)N\), using the contraction
\(E(1)\otimes_EV'(-1)\simeq V'\).  It satisfies
\(N(-1)\circ N=0:V'\to V'(-2)\).  If \(\widetilde N:V'\to V'\) is its
underlying linear operator after the chosen trivialization of \(E(1)\), and
\(F'\) is the action of \(\Phi_{0'}\), then
\[
 \widetilde N^2=0,\qquad
 F'\widetilde N(F')^{-1}=q^{-1}\widetilde N.
\]
Equivalently, \(N:V'\to V'(-1)\) is Frobenius-equivariant.
\end{theorem}

\begin{definition}[Type-3 monodromy filtration]
Put
\[
 B'=\operatorname{im}(\widetilde N),\qquad
 K'=\ker(\widetilde N),\qquad
 Q'=V'/K'.
\]
The length-two monodromy filtration is
\[
 0=M_{-2}V'\subset M_{-1}V'=B'
 \subset M_0V'=K'\subset M_1V'=V'.
\]
Its nonzero graded terms are
\[
 \operatorname{Gr}_{-1}^MV'=B',\qquad
 \operatorname{Gr}_{0}^MV'=K'/B',\qquad
 \operatorname{Gr}_{1}^MV'=Q'.
\]
\end{definition}

\begin{theorem}[Type-3 graded monodromy map and inverse]
The filtration is Frobenius-stable,
\(K'=(V')^{I_{0'}}\), and \(Q'=W_1(V')\).  The logarithm induces a
Frobenius-equivariant isomorphism
\[
 \overline N:\operatorname{Gr}_{1}^MV'=Q'
       \xrightarrow{\sim}\operatorname{Gr}_{-1}^MV'(-1)=B'(-1).
\]
Its inverse \(s_N:B'(-1)\to Q'\) satisfies
\[
 s_N\overline N=1_{Q'},\qquad
 \overline Ns_N=1_{B'(-1)}.
\]
Both identities commute with geometric Frobenius and with finite extension
of the coefficient and residue fields.
\end{theorem}

\begin{theorem}[Type-3 upper graded-weight bound]
If \(\mathcal F'_0\) is punctually \(\iota\)-pure of weight \(w'\), then
the Frobenius module \(Q'=\operatorname{Gr}_{1}^MV'\) has weights at most
\(w'+1\).
\end{theorem}

\begin{theorem}[Type-3 dual lower graded-weight bound]
For the dual monodromy filtration on \((V')^\vee\), evaluation identifies
\((Q')^\vee\) with the degree \(-1\) dual graded term.  Local duality makes
this module have weights at most \(-w'-1\).  Equivalently, \(Q'\) has
weights at least \(w'+1\).  The dual of \(\overline N\) and the dual of
\(s_N\) are mutually inverse Frobenius-equivariant graded maps.
\end{theorem}

\begin{theorem}[Type-3 length-two monodromy]
Let \(\mathcal F_0\) satisfy the hypotheses of the type-3 Fourier-sheaf
theorem, and suppose that its Fourier output \(\mathcal F'_0\) is punctually
\(\iota\)-pure of a real weight \(w'\).  For
\(V'=\mathcal F'_{\bar\eta_{0'}}\), the Frobenius module
\[
                         W_1(V')=V'/(V')^{I_{0'}}
\]
with its induced geometric-Frobenius operator through
\(W_{0'}/I_{0'}\) is \(\iota\)-pure of weight \(w'+1\).
\end{theorem}

\begin{theorem}[Type-3 weight comparison]
Assume that \(\mathcal F_0\) is punctually \(\iota\)-pure of a real weight
\(w\) and that the type-3 Fourier output \(\mathcal F'_0\) is punctually
\(\iota\)-pure of a real weight \(w'\).  Then
\[
                              w'=w+1.
\]
Moreover, the Frobenius modules in the local Fourier exact sequence satisfy
\[
 W_1(V')\simeq V_\infty(-1)\quad\text{and is pure of weight }w+2,
\]
whereas \((V')^{I_{0'}}\) has weights at most \(w+1\).
\end{theorem}

\section{Laumon's affine-line theorem}

\begin{theorem}[Affine-line estimate, real-weight form]
Let \(j_0:U_0\hookrightarrow\Aone_k\) be dense.  Let \(\mathcal F_0\) be an
arithmetically irreducible lisse Weil sheaf on \(U_0\), punctually
\(\iota\)-pure of a real weight \(w\).  Assume that:
\begin{enumerate}
\item the geometric local system has no constant simple subquotient;
\item it is unramified at infinity;
\item either \(\mathcal F_0\) has a boundary-real \(\iota\)-conjugate
companion, or the type-3 Fourier output \(\mathcal F'_0\) is known
independently to be punctually \(\iota\)-pure of some real weight.
\end{enumerate}
Then
\[
 H_c^i(\Aone_{\Fbar},j_*\mathcal F)=0\quad(i\ne1),
\]
the remaining group \(H_c^1(\Aone_{\Fbar},j_*\mathcal F)\) is a
finite-dimensional \(E\)-vector space, and every eigenvalue \(\alpha\) of
geometric Frobenius on it satisfies
\[
       |\iota(\alpha)|\le q^{(w+1)/2}.
\]
Equivalently, \(H_c^1(\Aone_{\Fbar},j_*\mathcal F)\) has weights at most
\(w+1\).  The statement is compatible with finite extension of \(k\),
finite extension of \(E\), and scalar extension to \(\Qlb\).
\end{theorem}

\begin{corollary}[Integral-weight form]
Under hypotheses (1) and (2) of the affine-line theorem, if
\(w\in\mathbf Z\), then
\[
 H_c^i(\Aone_{\Fbar},j_*\mathcal F)=0\quad(i\ne1)
\]
and \(H_c^1(\Aone_{\Fbar},j_*\mathcal F)\) has weights at most \(w+1\).
\end{corollary}

\begin{corollary}[All embeddings]
Suppose that \(\mathcal F_0\) is punctually pure of the same real weight
\(w\) for every embedding
\(\iota:\Qlb\hookrightarrow\mathbf C\).  Suppose also, for each
embedding, that a boundary-real conjugate companion exists or that the
Fourier output is independently known to be pure.  Then every Frobenius
eigenvalue \(\alpha\) on
\(H_c^1(\Aone_{\Fbar},j_*\mathcal F)\) satisfies
\[
 |\iota(\alpha)|\le q^{(w+1)/2}
 \quad\text{for every }\iota:\Qlb\hookrightarrow\mathbf C.
\]
Thus the cohomology group has the embedding-independent upper weight bound
\(w+1\).
\end{corollary}

\section{References}

The local transforms and their rank, Swan-conductor, and slope formulas are
Laumon's Theorem~2.4.3.  The type-3 monodromy sequence and affine-line
estimate occur in Sections~4.3--4.4 of the same work.  The realness criterion
and local weight statements are due to Deligne.

\begin{thebibliography}{9}
\bibitem{Laumon}
G.~Laumon,
\emph{Transformation de Fourier, constantes d'equations fonctionnelles et
conjecture de Weil},
Publications Mathematiques de l'IHES \textbf{65} (1987), 131--210.

\bibitem{Deligne}
P.~Deligne,
\emph{La conjecture de Weil II},
Publications Mathematiques de l'IHES \textbf{52} (1980), 137--252.

\bibitem{BBD}
A.~Beilinson, J.~Bernstein, and P.~Deligne,
\emph{Faisceaux pervers},
Ast\'erisque \textbf{100} (1982).
\end{thebibliography}

\end{document}
